\documentclass[11pt,a4paper]{article}
\usepackage[T1]{fontenc}
\usepackage[utf8]{inputenc}
\usepackage[french]{babel}
\frenchsetup{SmallCapsFigTabCaptions=false}
\addto\captionsfrench{\renewcommand{\tablename}{Tableau}}
\usepackage[locale=FR]{siunitx}
\newcommand{\pct}[1]{\num{#1}\,\%}
\newcommand{\NTheorem}{Théorème}
\newcommand{\NProposition}{Proposition}
\newcommand{\NLemma}{Lemme}
\newcommand{\NCorollary}{Corollaire}
\newcommand{\NDefinition}{Définition}
\newcommand{\NAssumption}{Hypothèse}
\newcommand{\NExample}{Exemple}
\newcommand{\NRemark}{Remarque}
\newcommand{\NInsight}{À retenir}
\newcommand{\NCaveat}{Mise en garde}
\newcommand{\shorttitle}{Repo implicite Totem : théorie, contrôle et détection}
\newcommand{\draftname}{Document de travail}
\newcommand{\lang}{fr}
% Shared packages for the English and French theory notes. The main files load babel and siunitx with the
% language options and define the language-dependent names (\NTheorem, \NInsight, ...) before this file;
% hyperref, cleveref and then theorems.tex come after it.
\usepackage{lmodern}
\usepackage{amsmath,amssymb,amsthm,mathtools,bm}
\usepackage[a4paper,margin=2.35cm,headheight=14pt]{geometry}
\usepackage{graphicx,booktabs,makecell,array,xcolor,caption,enumitem,microtype}
\usepackage{fancyhdr}
\usepackage[most]{tcolorbox}
\usepackage{algorithm}
\usepackage[noend]{algpseudocode}

\definecolor{ink}{HTML}{0B0B0B}
\definecolor{ink2}{HTML}{52514E}
\definecolor{accent}{HTML}{2A78D6}
\definecolor{accentwash}{HTML}{EEF4FC}
\definecolor{warm}{HTML}{EB6834}
\definecolor{warmwash}{HTML}{FDF1EB}
\definecolor{rule}{HTML}{C3C2B7}

\captionsetup{font=small,labelfont=bf,margin=0.5cm}
\setlist{itemsep=2pt,topsep=4pt}
\renewcommand{\arraystretch}{1.15}
\setlength{\parskip}{3pt}

\newtcolorbox{insight}{enhanced,breakable,colback=accentwash,colframe=accent,boxrule=0.5pt,arc=1.5pt,
  left=6pt,right=6pt,top=4pt,bottom=4pt,fonttitle=\bfseries\small,coltitle=white,title=\NInsight}
\newtcolorbox{caveat}{enhanced,breakable,colback=warmwash,colframe=warm,boxrule=0.5pt,arc=1.5pt,
  left=6pt,right=6pt,top=4pt,bottom=4pt,fonttitle=\bfseries\small,coltitle=white,title=\NCaveat}

% ------------------------------------------------------------- notation
\newcommand{\R}{\mathbb{R}}
\newcommand{\E}{\mathbb{E}}
\newcommand{\Var}{\operatorname{Var}}
\newcommand{\Cov}{\operatorname{Cov}}
\newcommand{\sd}{\operatorname{sd}}
\newcommand{\N}{\mathcal{N}}
\newcommand{\T}{\mathrm{T}}
\newcommand{\FO}{\mathrm{FO}}
\newcommand{\diag}{\operatorname{diag}}
\newcommand{\rank}{\operatorname{rank}}
\newcommand{\one}{\mathbf{1}}
\newcommand{\tb}{\tilde{t}}
\newcommand{\Aset}{\mathcal{A}}
\newcommand{\Bset}{\mathcal{B}}
\newcommand{\bA}{\overline{b}_{\Aset}}
\newcommand{\DA}{\overline{D}_{\Aset}}
\newcommand{\DB}{\overline{D}_{\Bset}}
\newcommand{\code}[1]{\texttt{#1}}

\pagestyle{fancy}
\fancyhf{}
\renewcommand{\headrulewidth}{0.4pt}
\fancyhead[L]{\small\color{ink2}\shorttitle}
\fancyhead[R]{\small\color{ink2}\draftname}
\fancyfoot[C]{\small\thepage}

\usepackage[colorlinks=true,linkcolor=accent,citecolor=accent,urlcolor=accent]{hyperref}
\usepackage[nameinlink,capitalise,french]{cleveref}
% Theorem environments: must come after amsthm, hyperref and cleveref.
% ------------------------------------------------------------ theorems
\theoremstyle{plain}
\newtheorem{theorem}{\NTheorem}[section]
\newtheorem{proposition}[theorem]{\NProposition}
\newtheorem{lemma}[theorem]{\NLemma}
\newtheorem{corollary}[theorem]{\NCorollary}
\theoremstyle{definition}
\newtheorem{definition}[theorem]{\NDefinition}
\newtheorem{assumption}[theorem]{\NAssumption}
\newtheorem{example}[theorem]{\NExample}
\theoremstyle{remark}
\newtheorem{remark}[theorem]{\NRemark}


\crefname{theorem}{théorème}{théorèmes}
\Crefname{theorem}{Théorème}{Théorèmes}
\crefname{proposition}{proposition}{propositions}
\Crefname{proposition}{Proposition}{Propositions}
\crefname{lemma}{lemme}{lemmes}
\Crefname{lemma}{Lemme}{Lemmes}
\crefname{corollary}{corollaire}{corollaires}
\Crefname{corollary}{Corollaire}{Corollaires}
\crefname{definition}{définition}{définitions}
\Crefname{definition}{Définition}{Définitions}
\crefname{assumption}{hypothèse}{hypothèses}
\Crefname{assumption}{Hypothèse}{Hypothèses}
\crefname{example}{exemple}{exemples}
\Crefname{example}{Exemple}{Exemples}
\crefname{remark}{remarque}{remarques}
\Crefname{remark}{Remarque}{Remarques}
\crefname{algorithm}{algorithme}{algorithmes}
\Crefname{algorithm}{Algorithme}{Algorithmes}
\crefname{section}{section}{sections}
\Crefname{section}{Section}{Sections}
\crefname{appendix}{annexe}{annexes}
\Crefname{appendix}{Annexe}{Annexes}
\crefname{figure}{figure}{figures}
\Crefname{figure}{Figure}{Figures}
\crefname{table}{tableau}{tableaux}
\Crefname{table}{Tableau}{Tableaux}
\crefname{equation}{équation}{équations}
\Crefname{equation}{Équation}{Équations}
\floatname{algorithm}{Algorithme}
\algrenewcommand\algorithmicrequire{\textbf{Entrées :}}
\algrenewcommand\algorithmicensure{\textbf{Sorties :}}
\algrenewcommand\algorithmicfor{\textbf{pour}}
\algrenewcommand\algorithmicdo{\textbf{faire}}
\algrenewcommand\algorithmicif{\textbf{si}}
\algrenewcommand\algorithmicthen{\textbf{alors}}
\algrenewcommand\algorithmicelse{\textbf{sinon}}
\algrenewcommand\algorithmicrepeat{\textbf{répéter}}
\algrenewcommand\algorithmicuntil{\textbf{jusqu'à}}
\algrenewcommand\algorithmicreturn{\textbf{renvoyer}}
\newcommand{\vNTest}{\num{5400}}
\newcommand{\vNTrain}{\num{300}}
\newcommand{\vNTestPerCat}{\num{150}}
\newcommand{\vNCat}{4}
\newcommand{\vNScen}{9}
\newcommand{\vFPR}{\pct{2}}
\newcommand{\vFPRmin}{\pct{2}}
\newcommand{\vFPRmax}{\pct{3}}
\newcommand{\vMeanNaive}{\pct{57}}
\newcommand{\vFwdNaive}{\pct{35}}
\newcommand{\vFoNaive}{\pct{60}}
\newcommand{\vPanelNaive}{\pct{42}}
\newcommand{\vBucketNaive}{\pct{29}}
\newcommand{\vTimingNaive}{\pct{55}}
\newcommand{\vSpotNaive}{\pct{89}}
\newcommand{\vSeasonNaive}{\pct{90}}
\newcommand{\vGrossNaive}{\pct{70}}
\newcommand{\vPowNaiveA}{\pct{2}}
\newcommand{\vPowNaiveB}{\pct{11}}
\newcommand{\vPowNaiveC}{\pct{27}}
\newcommand{\vPowNaiveD}{\pct{38}}
\newcommand{\vPowNaiveE}{\pct{48}}
\newcommand{\vPowNaiveF}{\pct{52}}
\newcommand{\vPowNaiveG}{\pct{50}}
\newcommand{\vPowNaiveH}{\pct{84}}
\newcommand{\vMeanShape}{\pct{64}}
\newcommand{\vFwdShape}{\pct{87}}
\newcommand{\vFoShape}{\pct{3}}
\newcommand{\vPanelShape}{\pct{24}}
\newcommand{\vBucketShape}{\pct{41}}
\newcommand{\vTimingShape}{\pct{74}}
\newcommand{\vSpotShape}{\pct{64}}
\newcommand{\vSeasonShape}{\pct{92}}
\newcommand{\vGrossShape}{\pct{71}}
\newcommand{\vPowShapeA}{\pct{2}}
\newcommand{\vPowShapeB}{\pct{11}}
\newcommand{\vPowShapeC}{\pct{24}}
\newcommand{\vPowShapeD}{\pct{42}}
\newcommand{\vPowShapeE}{\pct{64}}
\newcommand{\vPowShapeF}{\pct{69}}
\newcommand{\vPowShapeG}{\pct{72}}
\newcommand{\vPowShapeH}{\pct{89}}
\newcommand{\vMeanGmmCurve}{\pct{69}}
\newcommand{\vFwdGmmCurve}{\pct{74}}
\newcommand{\vFoGmmCurve}{\pct{2}}
\newcommand{\vPanelGmmCurve}{\pct{36}}
\newcommand{\vBucketGmmCurve}{\pct{47}}
\newcommand{\vTimingGmmCurve}{\pct{65}}
\newcommand{\vSpotGmmCurve}{\pct{96}}
\newcommand{\vSeasonGmmCurve}{\pct{94}}
\newcommand{\vGrossGmmCurve}{\pct{80}}
\newcommand{\vPowGmmCurveA}{\pct{3}}
\newcommand{\vPowGmmCurveB}{\pct{12}}
\newcommand{\vPowGmmCurveC}{\pct{25}}
\newcommand{\vPowGmmCurveD}{\pct{53}}
\newcommand{\vPowGmmCurveE}{\pct{67}}
\newcommand{\vPowGmmCurveF}{\pct{79}}
\newcommand{\vPowGmmCurveG}{\pct{79}}
\newcommand{\vPowGmmCurveH}{\pct{95}}
\newcommand{\vMeanGmmDiff}{\pct{73}}
\newcommand{\vFwdGmmDiff}{\pct{78}}
\newcommand{\vFoGmmDiff}{\pct{75}}
\newcommand{\vPanelGmmDiff}{\pct{39}}
\newcommand{\vBucketGmmDiff}{\pct{56}}
\newcommand{\vTimingGmmDiff}{\pct{69}}
\newcommand{\vSpotGmmDiff}{\pct{100}}
\newcommand{\vSeasonGmmDiff}{\pct{94}}
\newcommand{\vGrossGmmDiff}{\pct{86}}
\newcommand{\vPowGmmDiffA}{\pct{3}}
\newcommand{\vPowGmmDiffB}{\pct{13}}
\newcommand{\vPowGmmDiffC}{\pct{35}}
\newcommand{\vPowGmmDiffD}{\pct{67}}
\newcommand{\vPowGmmDiffE}{\pct{77}}
\newcommand{\vPowGmmDiffF}{\pct{84}}
\newcommand{\vPowGmmDiffG}{\pct{83}}
\newcommand{\vPowGmmDiffH}{\pct{97}}
\newcommand{\vMeanStruct}{\pct{73}}
\newcommand{\vFwdStruct}{\pct{97}}
\newcommand{\vFoStruct}{\pct{3}}
\newcommand{\vPanelStruct}{\pct{24}}
\newcommand{\vBucketStruct}{\pct{55}}
\newcommand{\vTimingStruct}{\pct{71}}
\newcommand{\vSpotStruct}{\pct{98}}
\newcommand{\vSeasonStruct}{\pct{94}}
\newcommand{\vGrossStruct}{\pct{88}}
\newcommand{\vPowStructA}{\pct{2}}
\newcommand{\vPowStructB}{\pct{5}}
\newcommand{\vPowStructC}{\pct{14}}
\newcommand{\vPowStructD}{\pct{44}}
\newcommand{\vPowStructE}{\pct{85}}
\newcommand{\vPowStructF}{\pct{100}}
\newcommand{\vPowStructG}{\pct{99}}
\newcommand{\vPowStructH}{\pct{99}}
\newcommand{\vAttr}{\pct{88}}
\newcommand{\vFoMismark}{\pct{80}}
\newcommand{\vFoClean}{\pct{2}}
\newcommand{\vFoMismarkEuStock}{\pct{97}}
\newcommand{\vKappaEuStock}{\num{1.60}}
\newcommand{\vLevelScaleEuStock}{\num{1.00}}
\newcommand{\vFoMismarkUsStock}{\pct{100}}
\newcommand{\vKappaUsStock}{\num{1.88}}
\newcommand{\vLevelScaleUsStock}{\num{1.05}}
\newcommand{\vFoMismarkHtb}{\pct{24}}
\newcommand{\vKappaHtb}{\num{0.96}}
\newcommand{\vLevelScaleHtb}{\num{1.00}}
\newcommand{\vFoMismarkIndex}{\pct{100}}
\newcommand{\vKappaIndex}{\num{3.58}}
\newcommand{\vLevelScaleIndex}{\num{1.18}}
\newcommand{\vZcrit}{\num{3.39}}
\newcommand{\vPowThFour}{\pct{73.0}}
\newcommand{\vPowThSix}{\pct{99.5}}
\newcommand{\vPowThThree}{\pct{54}}
\newcommand{\vNTemplates}{22}
\newcommand{\vGapOneY}{\num{6.5}}
\newcommand{\vGapFiveY}{\num{51.6}}
\newcommand{\vBookATotal}{\num{-1.74}}
\newcommand{\vMatA}{\pct{99.8}}
\newcommand{\vMatB}{\pct{60.8}}
\newcommand{\vMatC}{\pct{77.6}}
\newcommand{\vParF}{\num{102.30}}
\newcommand{\vParDF}{\num{0.9700}}
\newcommand{\vParFtwo}{\num{102.30}}
\newcommand{\vParDFtwo}{\num{0.9894}}
\newcommand{\vToyC}{\num{3.11}}
\newcommand{\vToyCsidak}{\num{3.19}}
\newcommand{\vToyQ}{\num{21.0}}
\newcommand{\vToyBeta}{\num{36.1}}
\newcommand{\vToyM}{7}
\newcommand{\vLevMin}{\num{0.32}}
\newcommand{\vLevMax}{\num{0.80}}
\newcommand{\vLevMean}{\num{0.545}}
\newcommand{\vNewtonA}{\num{3.1e-05}}
\newcommand{\vNewtonB}{\num{3.5e-12}}
\newcommand{\vJumpA}{\num{0.14}}
\newcommand{\vJumpB}{\num{0.24}}
\newcommand{\vJumpC}{\num{0.28}}
\newcommand{\vJumpD}{\num{0.39}}
\newcommand{\vRuntime}{\num{4}}


\title{\textbf{Repo implicite Totem}\\[4pt]\large Théorie, contrôle indépendant des dividendes et détection des courbes
aberrantes\\[2pt]\normalsize Note technique avec démonstrations, exemples détaillés et banc d'essai synthétique}
\author{Contrôle indépendant des prix -- Dérivés actions}
\date{14 octobre 2026 \quad\textbullet\quad Document de travail}

\begin{document}
\maketitle
\thispagestyle{empty}

\begin{abstract}
\noindent Les paramètres de repo utilisés pour le contrôle indépendant des prix sont implicités de Totem : les
forwards à partir des prix de consensus des calls et des puts, les dividendes à partir du service dividendes de Totem,
distinct. Comme ces deux services sont alimentés par des panels de contributeurs différents dont la taille varie avec
la maturité, la courbe de repo implicite hérite du bruit et des artefacts des deux, et présente parfois des formes que
le marché ne justifie pas. Cette note donne (i)~une décomposition exacte du repo implicite en repo du panel options
plus un \emph{écart de panel}, avec le bruit, les sauts et les biais que chaque mécanisme produit ; (ii)~un théorème de
non-identifiabilité qui établit qu'aucune méthode fondée sur les seules données de consensus ne peut supprimer cet
écart ; (iii)~une identité de contrôle qui démontre que la somme des lignes de contrôle dividende et repo ne dépend
pas du panel dividendes pour une exposition linéaire en forward, accompagnée d'une mesure de matérialité calculée à
partir des seules sensibilités ; (iv)~un détecteur structurel fondé sur des filtres adaptés et un test du rapport de
vraisemblance généralisé, calibré sur un historique sain, qui nomme le mécanisme à l'origine de chaque alerte et
répare la courbe ; (v)~un banc d'essai sur un marché Totem synthétique construit contributeur par contributeur,
qui le compare à des contrôles naïfs et à un regroupement par mélange gaussien de coefficients B-spline. Chaque
résultat est démontré, illustré par un exemple chiffré et vérifié dans le package Python qui accompagne la note.
\end{abstract}

\tableofcontents
\newpage

\section{Introduction}\label{sec-intro}

\subsection{Contexte}
Une fonction de contrôle indépendant des prix (IPV) contrôle les paramètres que le front office (FO) utilise pour
valoriser et gérer le risque des dérivés actions. Pour le repo, la référence n'est pas observée directement. Elle est
\emph{implicitée} : Totem publie des prix de consensus de calls et de puts, dont on extrait les forwards par la
parité call--put, et un service Totem distinct publie des dividendes de consensus. La courbe de repo qui réconcilie
les deux est ensuite comparée aux repos marqués par le FO, et les écarts sont multipliés par les sensibilités FO pour
produire des lignes de contrôle (écarts de valorisation, provisions, ajustements de valorisation prudente).

Deux caractéristiques de Totem rendent l'exercice délicat. Le service options et le service dividendes sont alimentés
par des contributeurs différents, et dans chaque service le nombre de contributeurs varie avec la maturité ; une action
typique peut compter quinze contributeurs à un an sur les options et dix sur les dividendes, moins au-delà de deux ans.
Le repo implicite mélange les deux moyennes de consensus : il hérite donc du bruit des deux et des artefacts de leur
décalage. Les contrôleurs le constatent sous forme de courbes de repo implicite irrégulières et de sauts injustifiés.

\subsection{Questions traitées}
Cette note répond à quatre questions :
\begin{enumerate}[label=(\alph*)]
\item \textbf{D'où viennent les courbes de repo implicite irrégulières ?} La \cref{sec-anatomy} décompose exactement
le repo implicite en repo du panel options plus un écart de panel, et établit la signature de chaque mécanisme
(écart de panel, sortie de contributeurs, date de détachement, profil de détachement périmé, erreur sur une période de
dividendes, décalage du spot, définition brute/nette).
\item \textbf{Une meilleure méthode statistique peut-elle retrouver le « vrai » repo ?} Non : le \cref{thm-nonid}
montre que les données de consensus ne permettent pas d'identifier la répartition entre repo et dividendes du panel
options.
\item \textbf{Faut-il une courbe de repo Totem propre pour contrôler le repo ?} Pour une exposition linéaire en
forward, non : les \cref{thm-control,thm-exact} démontrent que la somme des lignes dividende et repo ne dépend pas
du panel dividendes, et la \cref{prop-materiality} mesure, à partir des seules sensibilités, l'importance de la
répartition pour un portefeuille donné.
\item \textbf{Comment détecter les courbes aberrantes ?} Les \cref{sec-noise,sec-detection} construisent un
détecteur structurel : un modèle de bruit calculé à partir des propres statistiques de Totem, un gabarit de filtre
adapté par mécanisme, un taux de fausses alertes calibré et une explication en unités naturelles pour chaque alerte.
La \cref{sec-benchmark} le teste, et la \cref{sec-alternatives} explique la place des coefficients B-spline, des
mélanges gaussiens et des petits réseaux de neurones.
\end{enumerate}

\subsection{Principaux résultats}
\begin{insight}
\begin{itemize}[leftmargin=12pt]
\item Au premier ordre, $b_{\text{implicite}}=\bA+(\DA-\DB)/\bigl(T(S-D)\bigr)$ (\cref{prop-gap}). L'écart seul
produit \vGapOneY\,pb de bruit de repo à un an et \vGapFiveY\,pb à cinq ans pour une action typique.
\item Les mécanismes ont des signatures distinctes dans l'\emph{espace de portage} $B=b\,T$ : une constante (spot),
une marche (montant de dividende, changement de panel), un pic (forward aberrant, date de détachement), un motif
saisonnier (profil d'allocation).
\item Pour un portefeuille linéaire en forward, le P\&L Totem en réévaluation complète est \emph{exactement}
indépendant du panel dividendes lorsque les échéances sont des piliers Totem (\cref{thm-exact}) ; la décomposition en
lignes le reproduit au premier ordre (\cref{thm-control}). Le panel dividendes n'intervient qu'à travers
l'exposition non-forward $g=s_D+s_bJ$.
\item Les statistiques de test du détecteur sont des filtres adaptés exacts (\cref{thm-glrt}). Une seule échelle de
variance par catégorie, calibrée sur un historique sain, fixe exactement le taux de fausses alertes (\cref{prop-kappa}).
\item Sur le banc d'essai (\vNTest{} sous-jacents synthétiques de test, chaque détecteur à \vFPR{} de fausses
alertes), le détecteur structurel et un mélange gaussien sur les coefficients B-spline de l'écart Totem--FO détectent
en moyenne une part comparable des anomalies (\vMeanStruct{} et \vMeanGmmDiff). Le détecteur structurel trouve en
outre les forwards aberrants (\vFwdStruct{} contre au plus \vFwdShape), explique correctement \vAttr{} de ses alertes
et ne confond pas une véritable erreur de marquage FO avec une courbe Totem aberrante (\vFoStruct{} contre
\vFoGmmDiff).
\end{itemize}
\end{insight}

\subsection{Guide de lecture}
La \cref{sec-model} fixe les notations et démontre les résultats élémentaires sur lesquels repose la suite
(bootstrap, parité, allocation). La \cref{sec-anatomy} pose le diagnostic ; la \cref{sec-control} établit l'identité
de contrôle ; les \cref{sec-noise,sec-detection} construisent le détecteur ; la \cref{sec-benchmark} apporte les
éléments empiriques ; la \cref{sec-alternatives} compare avec d'autres approches ; la \cref{sec-ops} décrit la mise en
œuvre. Les démonstrations suivent chaque énoncé ; les calculs plus longs figurent en \cref{app-proofs}. Chaque exemple
chiffré est produit par le code du dépôt (\cref{app-repro}).

\section{Modèle et notations}\label{sec-model}

\subsection{Forwards avec dividendes discrets}
Le temps est mesuré en années à partir de la date de valorisation. Soit $S$ le spot, $r(t)$ le taux zéro-coupon en
composition continue utilisé pour l'actualisation et le financement, et $b(t)$ la courbe zéro-coupon de \emph{repo}
(coût d'emprunt du titre), de sorte que détenir l'action financée par un repo rapporte $r-b$. On note $R(t)=r(t)\,t$,
$B(t)=b(t)\,t$ et $G(t)=e^{R(t)-B(t)}$. Des dividendes en numéraire $d_i$ sont détachés aux dates $t_i$.

\begin{definition}[Forward]\label{def-forward}
Le forward de livraison $T$ vaut
\begin{equation}\label{eq-forward}
F(T)=G(T)\bigl(S-D(T)\bigr),\qquad D(T)=\sum_{t_i\le T}\frac{d_i}{G(t_i)}.
\end{equation}
$D(T)$ est le \emph{dividende cumulé actualisé au portage} : les dividendes sont actualisés au taux de portage $r-b$.
\end{definition}

\Cref{eq-forward} est l'argument de réplication habituel : acheter l'action financée à $r-b$ et reporter les
dividendes. Avec des courbes plates, $F(T)=Se^{(r-b)T}-\sum_{t_i\le T}d_ie^{(r-b)(T-t_i)}$. On note que $D$ dépend de la
courbe de repo à travers $G(t_i)$ ; ce couplage est faible (quelques points de base de repo par an sur l'actualisation
des dividendes), mais le bootstrap ci-dessous le traite exactement.

\subsection{L'espace de portage}
\begin{definition}[Portage]\label{def-carry}
Le \emph{repo cumulé} ou \emph{portage} est $B(T)=b(T)\,T$. Pour une grandeur observée aux piliers $T_1<\dots<T_K$, on
note $B_k=B(T_k)$, $b_k=b(T_k)$, $D_k=D(T_k)$, $F_k=F(T_k)$.
\end{definition}

\begin{lemma}[Additivité en portage]\label{lem-additivity}
$\ln F(T)=R(T)-B(T)+\ln\bigl(S-D(T)\bigr)$. Par conséquent, des perturbations du forward, des dividendes et du portage
vérifient, au premier ordre,
\begin{equation}\label{eq-additivity}
\delta B(T)=-\,\delta\ln F(T)-\frac{\delta D(T)}{S-D(T)}+\delta R(T).
\end{equation}
\end{lemma}
\begin{proof}
On prend le logarithme de \cref{eq-forward} et on différentie, avec $\delta\ln(S-D)=-\delta D/(S-D)+O(\delta D^2)$.
\end{proof}

\begin{insight}
Dans l'espace de portage, les mécanismes s'additionnent et rien n'est amplifié. Une erreur de portage $\varepsilon$
à la maturité $T$ est une erreur de repo $\varepsilon/T$ : la même erreur de $0{,}1\,\%$ sur le forward vaut $120$\,pb de
repo à un mois et $2$\,pb à cinq ans. C'est pourquoi les courbes de repo implicite paraissent si bruitées sur le court
terme, et pourquoi tous les tests de cette note portent sur $B$ et non sur $b$.
\end{insight}

\subsection{L'application repo implicite}
Étant donnés des forwards $F_1,\dots,F_K$ aux piliers, la courbe des taux et les dividendes, le \emph{repo implicite}
est la courbe de repo, affine par morceaux en $B$ (repo forward constant entre piliers, $B(0)=0$), qui reprice les
forwards. On la résout pilier par pilier. On fixe $B_1,\dots,B_{k-1}$ et on pose $x=B_k$. Pour un dividende situé entre
le pilier précédent et $T_k$, $B(t_i)=B_{k-1}+w_i(x-B_{k-1})$ avec $w_i=(t_i-T_{k-1})/(T_k-T_{k-1})\in(0,1]$, de sorte que
le forward du modèle en $T_k$ vaut
\begin{equation}\label{eq-phik}
\Phi_k(x)=Se^{R_k-x}-\sum_{t_i\le T_{k-1}}d_i\,e^{(R_k-R_i)-(x-B(t_i))}
-\sum_{T_{k-1}<t_i\le T_k}d_i\,e^{(R_k-R_i)-(1-w_i)(x-B_{k-1})}.
\end{equation}

\begin{lemma}[Monotonie]\label{lem-monotone}
$\Phi_k'(x)\le-\Phi_k(x)$ pour tout $x$. En particulier $\Phi_k$ est strictement décroissante là où le forward est
positif, et $g_k(x)=\ln\Phi_k(x)-\ln F_k$ vérifie $g_k'(x)\le-1$.
\end{lemma}
\begin{proof}
En dérivant \cref{eq-phik} terme à terme,
$\Phi_k'(x)=-Se^{R_k-x}+\sum_{t_i\le T_{k-1}}d_ie^{(\cdot)}+\sum_{T_{k-1}<t_i\le T_k}(1-w_i)\,d_ie^{(\cdot)}$.
Chaque terme de dividende est positif et $1-w_i\le1$ ; remplacer $(1-w_i)$ par $1$ ne peut qu'augmenter le membre de
droite, qui vaut alors $-\Phi_k(x)$. En divisant par $\Phi_k>0$, on obtient $g_k'\le-1$.
\end{proof}

\begin{proposition}[Bootstrap bien posé]\label{prop-bootstrap}
On suppose $S-D(T_{k-1})>0$ pour les dividendes déjà résolus. Alors pour tout $F_k>0$ il existe un unique $B_k$ tel que
$\Phi_k(B_k)=F_k$ et $\Phi_k>0$ sur la branche de la solution. La méthode de Newton appliquée à $g_k$ converge
localement de façon quadratique, et globalement si on la sécurise par dichotomie.
\end{proposition}
\begin{proof}
Quand $x\to-\infty$,
\[
e^{x}\Phi_k(x)\to e^{R_k}\Bigl(S-\sum_{t_i\le T_{k-1}}d_ie^{-R_i+B(t_i)}\Bigr)=e^{R_k}\bigl(S-D(T_{k-1})\bigr)>0,
\]
car les termes de la période courante ne croissent que comme $e^{-(1-w_i)x}$ ; donc $\Phi_k\to+\infty$. Quand $x$ croît,
$\Phi_k$ décroît strictement tant qu'elle est positive (\cref{lem-monotone}) et soit tend vers $0$, soit change de
signe ; par continuité, elle prend chaque valeur de $(0,\infty)$ exactement une fois sur la branche où elle est
positive. Sur cette branche, $g_k$ est $C^2$ avec $|g_k'|\ge1$, donc l'itération de Newton $x\mapsto x-g_k(x)/g_k'(x)$
converge quadratiquement au voisinage de la racine ; la sécurité par dichotomie (conserver un encadrement $[\ell,u]$
avec $g_k(\ell)>0>g_k(u)$ et revenir à son milieu si nécessaire) rend la convergence globale.
\end{proof}

\begin{example}[Un bootstrap à la main]\label{ex-bootstrap}
Spot $100$, taux plat de $3\,\%$, dividendes de $1{,}5$ en $t=0{,}30$ et de $2{,}0$ en $t=0{,}80$, piliers $0{,}25$,
$0{,}5$ et $1$, repo vrai de $80$, $70$ et $60$\,pb. Les forwards et le repo implicite sont
\begin{center}\input{generated/\lang/tab_boot}\end{center}
Au troisième pilier, le résidu de Newton $|g_3|$ passe de \vNewtonA{} à \vNewtonB{} en une itération, puis à zéro à
la suivante. Sans dividende, $g_k$ est affine et une seule itération suffit.
\end{example}

\subsection{Forwards à partir des prix d'options de consensus}
\begin{lemma}[Parité et forward de consensus]\label{lem-parity}
Supposons que chaque contributeur $j$ cote des calls et des puts européens avec un facteur d'actualisation commun $DF$.
Alors
\begin{enumerate}[label=(\roman*)]
\item $C_j(K)-P_j(K)=DF\,(F_j-K)$ pour tout strike ;
\item si le consensus moyenne les mêmes contributeurs à chaque strike,
$\overline{C}(K)-\overline{P}(K)=DF\,(\overline{F}-K)$, où $\overline{F}$ est la moyenne des forwards des contributeurs ;
\item le strike $K^\ast$ où $\overline{C}=\overline{P}$ est égal à $\overline{F}$ et ne dépend pas de $DF$ ;
multiplier tous les prix par un même facteur ne déplace pas $K^\ast$.
\end{enumerate}
\end{lemma}
\begin{proof}
(i) est la parité call--put~\cite{hull}. (ii) s'obtient en moyennant (i), qui est linéaire. (iii) On annule (ii) ;
multiplier les deux prix par un facteur multiplie la droite $K\mapsto\overline{C}-\overline{P}$ sans déplacer sa racine.
\end{proof}

\begin{example}\label{ex-parity}
Avec des strikes $90$, $100$ et $110$, des puts à $2{,}1$, $5{,}4$ et $10{,}9$ et des calls construits à partir de
$F=102{,}3$ et $DF=0{,}97$, la régression pondérée de $\overline{C}-\overline{P}$ sur $K$ donne $F=\vParF$ et
$DF=\vParDF$. Multiplier tous les prix par $1{,}02$ (un facteur d'actualisation erroné) donne $F=\vParFtwo$ et
$DF=\vParDFtwo$ : le forward est inchangé.
\end{example}

\begin{remark}
Deux situations mettent (ii)--(iii) en défaut : un filtrage des valeurs aberrantes appliqué séparément par strike et
par type d'option (les contributeurs moyennés diffèrent alors entre calls et puts), et l'exercice américain (la parité
devient une inégalité). Le $DF$ ajusté, comparé à la courbe, donne un spread de financement implicite qui signale ces
échéances. Le contrôleur conserve sa propre courbe des taux des deux côtés de la comparaison, de sorte que les taux
s'éliminent dans l'écart Totem--FO.
\end{remark}

\subsection{Périodes de dividendes et allocation}
\begin{definition}[Allocation]\label{def-allocation}
Les niveaux de dividendes Totem $L_y$ sont des sommes de dividendes détachés dans la période $y=(e_{y-1},e_y]$
(en général des années civiles). Un profil attribue des poids $\pi_i\ge0$ avec $\sum_{i\in y}\pi_i=1$ aux dates de
détachement de chaque période ; les dividendes alloués sont $d_i=L_y\pi_i$. Le contrôleur prend normalement le
calendrier de détachement du FO comme profil.
\end{definition}

\begin{lemma}\label{lem-allocation}
À $G$ fixé, $D_k=\sum_yA_{ky}L_y$ avec $A_{ky}=\sum_{i\in y,\,t_i\le T_k}\pi_i/G(t_i)$.
\end{lemma}
\begin{proof}
On substitue $d_i=L_y\pi_i$ dans \cref{eq-forward} et on regroupe par période.
\end{proof}
$A_{ky}$ est la part, actualisée au portage, de la période $y$ détachée avant l'échéance $k$. Elle vaut $0$ ou
environ $1/G$, sauf pour les périodes qui chevauchent une échéance : c'est exactement là que le profil compte.

\subsection{Modèle des données Totem}
\begin{assumption}[Contributeurs cohérents]\label{ass-consistent}
Le contributeur $j$ a des vues $b_j(\cdot)$ et $D_j(\cdot)$ et cote des prix d'options cohérents avec
$F_j(T)=G_j(T)(S-D_j(T))$, avec la courbe des taux et le spot communs.
\end{assumption}
\begin{assumption}[Panels]\label{ass-panels}
À l'échéance $T_k$, le panel options $\Aset_k$ compte $n_{A,k}$ contributeurs ; pour la période $y$, le panel
dividendes $\Bset_y$ en compte $n_{B,y}$. Les panels sont emboîtés en maturité ($\Aset_{k+1}\subseteq\Aset_k$,
$\Bset_{y+1}\subseteq\Bset_y$). Les erreurs de vue sont indépendantes d'un contributeur à l'autre. Les deux services ont
$m$ contributeurs en commun ($m=0$ pour des panels disjoints).
\end{assumption}
\begin{assumption}[Consensus]\label{ass-consensus}
Le consensus est la moyenne arithmétique sur le panel. Totem publie, par échéance et par période, le consensus, la
dispersion entre contributeurs et leur nombre.
\end{assumption}

\begin{definition}[Repo implicite Totem et écart de portage]\label{def-implied}
Le repo implicite Totem $b^\T$ est le bootstrap des forwards de consensus $\overline{F}_k$ avec les dividendes de
consensus alloués $\overline{L}_y$. Avec le repo FO $b^\FO$ aux mêmes piliers, l'\emph{écart de portage} est
\begin{equation}\label{eq-y}
y_k=B^\T_k-B^\FO_k=\bigl(b^\T_k-b^\FO_k\bigr)\,T_k,\qquad k=1,\dots,K.
\end{equation}
\end{definition}

\section{Anatomie du repo implicite}\label{sec-anatomy}

\subsection{L'écart de panel}
Pour un pilier $T$, on note $\bA(T)$ et $\DA(T)$ les moyennes des vues de repo et de dividendes sur le panel options,
et $\DB(T)$ le consensus dividendes du panel dividendes après allocation. Les vues sont des perturbations d'un centre
commun, $b_j=b_0+\delta b_j$ et $D_j=D_0+\delta D_j$, d'ordre $\delta$.

\begin{proposition}[Décomposition par écart de panel]\label{prop-gap}
Sous les \cref{ass-consistent,ass-panels,ass-consensus},
\begin{equation}\label{eq-gap}
b^\T(T)=\bA(T)+\frac{\DA(T)-\DB(T)}{T\,\bigl(S-D_0(T)\bigr)}+O(\delta^2).
\end{equation}
\end{proposition}
\begin{proof}
Pour $j\in\Aset$, le \cref{lem-additivity} donne $\ln F_j=\ln F_0-T\,\delta b_j-\delta D_j/(S-D_0)+O(\delta^2)$. Le
forward de consensus est la moyenne des $F_j$, et le logarithme d'une moyenne égale la moyenne des logarithmes à un
terme de variance d'ordre $\delta^2$ près, d'où
$\ln\overline{F}=\ln F_0-T\,\overline{\delta b}_{\Aset}-\overline{\delta D}_{\Aset}/(S-D_0)+O(\delta^2)$. Le contrôleur
alloue le consensus dividendes, $\DB=D_0+\overline{\delta D}_{\Bset}$, et applique le bootstrap :
$B^\T=R-\ln\overline{F}+\ln(S-\DB)
=B_0+T\,\overline{\delta b}_{\Aset}+\bigl(\overline{\delta D}_{\Aset}-\overline{\delta D}_{\Bset}\bigr)/(S-D_0)+O(\delta^2)$.
On divise par $T$ et on utilise $b_0+\overline{\delta b}_{\Aset}=\bA$.
\end{proof}

Le repo implicite est le repo propre au panel options plus un \emph{écart de panel} : la différence entre la vue
dividendes des contributeurs options, jamais observée, et celle du panel dividendes, divisée par $T(S-D)$. Pour un
contributeur donné, les erreurs sur le forward et sur les dividendes se compensent dans le repo implicite ; entre deux
panels, elles ne se compensent pas.

\begin{corollary}[Taille de l'écart]\label{cor-gap}
Supposons les erreurs de vue dividendes en $T$ i.i.d.\ de variance $\sigma_D^2(T)$ dans les deux populations, et
supposons que les deux panels partagent $m$ contributeurs qui utilisent la même vue dividendes dans les deux services.
Alors
\begin{equation}\label{eq-gap-sd}
\sd\bigl[b^\T-\bA\bigr]=\frac{\sigma_D(T)}{T\,(S-D(T))}\sqrt{\frac1{n_A}+\frac1{n_B}-\frac{2m}{n_An_B}}.
\end{equation}
\end{corollary}
\begin{proof}
$\DA-\DB=\sum_j(a_j-c_j)\,\delta D_j$ avec $a_j=1/n_A$ sur $\Aset$, $c_j=1/n_B$ sur $\Bset$ et zéro ailleurs. Sa variance
vaut $\sigma_D^2\sum_j(a_j-c_j)^2=\sigma_D^2\bigl[m(1/n_A-1/n_B)^2+(n_A-m)/n_A^2+(n_B-m)/n_B^2\bigr]$, qui se simplifie
en $\sigma_D^2\,(1/n_A+1/n_B-2m/(n_An_B))$.
\end{proof}

Avec des panels identiques ($m=n_A=n_B$), l'écart disparaît ; avec des panels disjoints, il est maximal. La
dispersion $\sigma_D(T)$ d'un dividende cumulé croît avec $T$, à peu près comme $T$ ou plus vite : l'écart est donc
faible en portage sur le court terme mais ne diminue pas en repo sur le long terme.

\begin{example}[Taille de l'écart pour une action]\label{ex-gap}
Spot $100$, rendement en dividendes de $4\,\%$, dispersion des contributeurs de $1\,\%+3\,\%\cdot T$ du dividende
cumulé, tailles de panel de la catégorie synthétique « action UE ». La formule \eqref{eq-gap-sd} avec $m=0$ et une
simulation contributeur par contributeur sans désaccord sur le repo donnent
\begin{center}\input{generated/\lang/tab_gap}\end{center}
À un an, si six contributeurs étaient communs aux deux services ($m=6$, $n_A=15$, $n_B=12$), l'écart passerait de
$6{,}5$ à $4{,}8$\,pb.
\end{example}

\begin{figure}[ht]
\centering\includegraphics[width=\linewidth]{figures/\lang/panel_gap.pdf}
\caption{(a) Bruit de repo dû au seul écart de panel, \cref{eq-gap-sd}, pour trois configurations de panels. (b) Un
décalage de spot $\varepsilon$ est constant en portage et se lit donc $\varepsilon/T$ en repo (\cref{prop-spot}).}
\label{fig-gap}
\end{figure}

\subsection{Sortie de contributeurs}
\begin{proposition}[Moyennes de panels emboîtés]\label{prop-nested}
Soit $e_{jk}$ l'erreur du contributeur $j$ à la maturité $k$, indépendante d'un contributeur à l'autre, avec
$\Var e_{jk}=\sigma_k^2$ et $\Cov(e_{jk},e_{jl})=\rho_{kl}\sigma_k\sigma_l$. Si les panels sont emboîtés, de tailles
$n_k$, les moyennes $\bar e_k$ vérifient
\begin{equation}\label{eq-nested}
\Cov(\bar e_k,\bar e_l)=\frac{\rho_{kl}\,\sigma_k\sigma_l}{\max(n_k,n_l)}.
\end{equation}
En particulier, avec des vues persistantes ($\rho=1$, $\sigma_k=\sigma_l=\sigma$) et $n_l\le n_k$,
$\sd(\bar e_l-\bar e_k)=\sigma\sqrt{1/n_l-1/n_k}$.
\end{proposition}
\begin{proof}
$\Cov(\bar e_k,\bar e_l)=\frac{1}{n_kn_l}\sum_{j\le\min(n_k,n_l)}\rho_{kl}\sigma_k\sigma_l
=\frac{\min(n_k,n_l)}{n_kn_l}\rho_{kl}\sigma_k\sigma_l$, et $\min/(n_kn_l)=1/\max$. Ensuite
$\Var(\bar e_l-\bar e_k)=\sigma^2(1/n_l+1/n_k-2/n_k)=\sigma^2(1/n_l-1/n_k)$.
\end{proof}

\begin{example}
En unités de dispersion des contributeurs, le consensus saute de lui-même de $\vJumpA\,\sigma$ quand un panel passe
de $16$ à $12$ contributeurs, de $\vJumpB\,\sigma$ de $12$ à $7$, de $\vJumpC\,\sigma$ de $15$ à $7$ et de
$\vJumpD\,\sigma$ de $10$ à $4$. Un saut de cette taille lors d'une sortie de contributeurs est du bruit, pas une
information ; le modèle de bruit de la \cref{sec-noise} le reproduit grâce à \cref{eq-nested}.
\end{example}

\subsection{Désaccords sur la date de détachement}
\begin{proposition}[Désaccord de date]\label{prop-timing}
Soit un dividende $d_i$ détaché en $t_i\in(T_k,T_{k+1})$ dans l'allocation des dividendes, et supposons qu'une fraction
$p$ du panel options le place avant $T_k$ (tous le placent avant $T_{k+1}$). En notant $\tilde d_i=d_i/G(t_i)$, au premier
ordre :
\begin{enumerate}[label=(\roman*)]
\item l'écart de portage ne bouge qu'en $T_k$ : $\delta y_k=p\,\tilde d_i/(S-D_k)$, soit un pic de repo
$p\,\tilde d_i/\bigl(T_k(S-D_k)\bigr)$ ;
\item en repo forward entre piliers, c'est un dipôle : $+\delta y_k/(T_k-T_{k-1})$ puis $-\delta y_k/(T_{k+1}-T_k)$ ;
\item la variance des log-forwards des contributeurs en $T_k$ augmente de $p(1-p)\,\tilde d_i^{\,2}/(S-D_k)^2$.
\end{enumerate}
\end{proposition}
\begin{proof}
(i) En $T_k$, une fraction $p$ du panel options inclut le dividende, donc $\ln\overline{F}_k$ baisse de
$p\,\tilde d_i/(S-D_k)$ (\cref{lem-additivity}) tandis que le $D^\T_k$ alloué l'exclut ; par \cref{eq-additivity},
$B^\T_k$ augmente d'autant. Avant $T_k$ personne n'inclut le dividende, après $T_{k+1}$ tout le monde l'inclut :
aucun autre pilier ne bouge. (ii) Le repo forward est la pente de $B$ entre piliers. (iii) Formule de la variance
totale : les contributeurs forment deux groupes dont les log-forwards moyens diffèrent de $\tilde d_i/(S-D_k)$, de poids
$p$ et $1-p$.
\end{proof}

\begin{example}\label{ex-timing}
Dans l'exemple de la \cref{fig-gallery}(c), un dividende de $2{,}18$ est détaché $17$ jours après l'échéance à six mois,
et $37\,\%$ du panel options le placent avant. Le pic vaut $0{,}37\times2{,}18/(0{,}5\times98)\approx1{,}6\,\%$, soit
environ $165$\,pb de repo sur un seul pilier. La dispersion des contributeurs à cette échéance est presque dix fois
celle de ses voisines, conformément à (iii).
\end{example}

\subsection{Profil de détachement périmé}
\begin{proposition}[Décalage de profil]\label{prop-profile}
Supposons que les contributeurs options utilisent les vraies dates $t_i$, tandis que le profil d'allocation les
place en $t_i+\Delta$. Alors, pour $\Delta>0$ (symétriquement pour $\Delta<0$),
\[
\delta y_k=\sum_{i:\,T_k\in[t_i,\,t_i+\Delta)}\frac{\tilde d_i}{S-D_k}\,:
\]
seuls bougent les piliers situés entre une vraie date de détachement et sa copie décalée.
\end{proposition}
\begin{proof}
$D^\T_k-D_k=\sum_i\tilde d_i\bigl(\one\{t_i+\Delta\le T_k\}-\one\{t_i\le T_k\}\bigr)$, qui vaut $-\tilde d_i$ exactement
lorsque $t_i\le T_k<t_i+\Delta$ ; on applique \cref{eq-additivity}.
\end{proof}
Pour les payeurs saisonniers, les piliers touchés reviennent chaque année (par exemple les échéances d'avril des
actions européennes), ce qui produit le zigzag que les contrôleurs reconnaissent. Les fins de période ne sont jamais
touchées, car les totaux par période sont justes.

\subsection{Erreurs sur une période de dividendes}
\begin{proposition}[Erreur de période]\label{prop-bucket}
Si le niveau Totem de la période $y$ est erroné de $\delta L$, alors $\delta y_k=-A_{ky}\,\delta L/(S-D_k)$. Pour une
période à une seule date de détachement, c'est une marche qui commence au premier pilier après le détachement ; en
repo, c'est un saut qui décroît comme $1/T$.
\end{proposition}
\begin{proof}
Le \cref{lem-allocation} donne $\delta D_k=A_{ky}\delta L$ ; on applique \cref{eq-additivity} à forwards inchangés.
$A_{ky}$ est nul avant le détachement et constant ($\approx1/G(t_i)$) après.
\end{proof}

\subsection{Décalage du spot et dividendes proches}
\begin{proposition}[Décalage du spot]\label{prop-spot}
Si les contributeurs options valorisent avec un spot $S(1+\varepsilon)$ alors que le contrôleur utilise $S$,
\begin{equation}\label{eq-spot}
y_k=-\ln\Bigl(1+\frac{\varepsilon S}{S-D_k}\Bigr)\approx-\frac{\varepsilon S}{S-D_k}\qquad\text{pour tout }k,
\end{equation}
quasi constant d'un pilier à l'autre (exactement constant sans dividende), donc une erreur de repo proche de
$-\varepsilon/T_k$. Un dividende antérieur au premier pilier, erroné de $\delta d$, produit la même signature,
$-\delta\tilde d/(S-D_k)$.
\end{proposition}
\begin{proof}
Le forward de chaque contributeur devient $G(S(1+\varepsilon)-D)$, donc $\ln\overline{F}_k$ varie de
$\ln\bigl((S(1+\varepsilon)-D_k)/(S-D_k)\bigr)$ ; par \cref{eq-additivity}, $B^\T_k$ varie de l'opposé. Pour le
dividende proche, $D_k$ varie de $\delta\tilde d$ à tous les piliers.
\end{proof}

\subsection{Définitions brut/net et biais proportionnels}
\begin{proposition}[Biais proportionnel sur les dividendes]\label{prop-gross}
Si chaque dividende Totem vaut $(1+\tau)$ fois le dividende contenu dans les forwards (par exemple brut contre net de
retenue à la source), alors $y_k\approx-\tau D_k/(S-D_k)$. Pour un rendement continu $q$, $D_k\approx SqT_k$ et
$y_k\approx-\tau q\,T_k$, qui a exactement la forme d'un écart de niveau de repo de $-\tau q$.
\end{proposition}
\begin{proof}
$\delta D_k=\tau D_k$ dans \cref{eq-additivity} ; pour un rendement, $D_k=S(1-e^{-qT_k})\approx SqT_k$ et $S-D_k\approx S$.
\end{proof}

\subsection{Ce que les données ne peuvent pas dire}
\begin{theorem}[Non-identifiabilité de la répartition]\label{thm-nonid}
Supposons que les prix d'options européennes ne dépendent des vues d'un contributeur qu'à travers ses forwards. Pour
toute transformation au niveau des contributeurs $(b_j,D_j)\mapsto(b_j+\delta_j,D'_j)$ avec
$\delta_j(T)\,T=\ln(S-D'_j(T))-\ln(S-D_j(T))$, tous les prix d'options, et donc toutes les statistiques du panel options
(consensus, dispersion, nombre), sont inchangés ; les statistiques du panel dividendes le sont aussi. Le repo implicite
$b^\T$, fonction de ces statistiques, est donc inchangé, alors que le repo du panel options $\bA$ varie de la moyenne
des $\delta_j$. Par conséquent, l'écart de panel $b^\T-\bA$ n'est pas une fonction des données observables.
\end{theorem}
\begin{proof}
La transformation laisse $\ln F_j=R-B_j+\ln(S-D_j)$ inchangé pour tout $j$ et tout $T$ ; les prix d'options de chaque
contributeur sont donc inchangés, ainsi que toute statistique calculée à partir d'eux. Le panel dividendes n'est pas
touché. $b^\T$ n'est calculé qu'à partir de ces statistiques, il est donc inchangé, alors que $\bA$ se déplace de la
moyenne des $\delta_j$, qui peut prendre n'importe quelle valeur.
\end{proof}

\begin{corollary}
Aucun estimateur construit à partir des seules données de consensus Totem (lissage de courbe, regroupement, réseaux de
neurones) ne peut supprimer l'écart de panel. Il reste possible d'en modéliser la taille (\cref{sec-noise}), de
reconnaître les mécanismes structurés à leur signature (\cref{sec-detection}) et de se passer de la répartition là où
l'exposition ne l'exige pas (\cref{sec-control}).
\end{corollary}

\begin{remark}[Ce qui lève l'indétermination]
Les options américaines portent un peu d'information sur la répartition à travers l'exercice anticipé ; des
contributeurs communs réduisent l'écart (\cref{cor-gap}) ; des ancres externes fixent directement l'un des deux
côtés : contrats à terme cotés sur dividendes, contrats à terme de rendement total qui cotent un spread de financement,
commissions de prêt-emprunt de titres pour le court terme des actions.
\end{remark}

\begin{insight}
Résumé des signatures dans l'espace de portage (gabarits du détecteur entre crochets) :
\begin{center}\small
\begin{tabular}{p{4.3cm}p{4.9cm}p{4.4cm}}
\toprule
Mécanisme & Écart de portage $y$ & Écart de repo\\
\midrule
Écart de panel (\ref{prop-gap}) & bruit lisse plus incréments aux détachements & bruit à peu près plat en maturité\\
Sortie de contributeurs (\ref{prop-nested}) & marche à la sortie [marche] & saut décroissant en $1/T$\\
Désaccord de date (\ref{prop-timing}) & un pilier [pic] & pic, dipôle en repo forward\\
Profil périmé (\ref{prop-profile}) & piliers en saison [saison, pic] & zigzag\\
Erreur de période (\ref{prop-bucket}) & marche après le détachement [marche] & saut décroissant en $1/T$\\
Décalage du spot (\ref{prop-spot}) & constante [niveau] & hyperbole $\varepsilon/T$\\
Biais proportionnel (\ref{prop-gross}) & rampe de marches ; $\approx\beta T$ pour un rendement & marches ; un niveau pour un indice\\
Erreur de marquage FO & $\delta b(T)\,T$ lisse [modèle nul] & niveau ou pente lisse\\
\bottomrule
\end{tabular}
\end{center}
\end{insight}

\section{Contrôle du repo à partir des contrôles forward et dividendes}\label{sec-control}

\subsection{Lignes de contrôle}
On regroupe les taux de repo zéro-coupon aux piliers et les niveaux de dividendes par période dans
$\theta=(b,L)\in\R^K\times\R^Y$. Le spot, la courbe des taux et le profil d'allocation sont fixés, de sorte que les
forwards aux piliers forment une application régulière $F(\theta)\in\R^K$. D'après la \cref{prop-bootstrap},
l'\emph{application repo implicite} $\beta(F,L)$, l'unique $b$ tel que $F(b,L)=F$, est bien définie.

Les marquages FO sont $\theta^\FO=(b^\FO,L^\FO)$. Totem fournit des forwards $F^\T$ et des niveaux de dividendes $L^\T$,
d'où $b^\T=\beta(F^\T,L^\T)$. Une position a une valeur $V(\theta)$, avec des sensibilités FO
$s_b=\partial_bV(\theta^\FO)\in\R^K$ et $s_D=\partial_LV(\theta^\FO)\in\R^Y$. On note $\Delta L=L^\T-L^\FO$,
$\Delta b=b^\T-b^\FO$ et $\Delta b^F=\beta(F^\T,L^\FO)-b^\FO$ : l'écart de repo implicité par les forwards Totem
\emph{avec les dividendes FO}.

\begin{definition}[Lignes de contrôle]\label{def-lines}
\[
\text{ligne dividende}=s_D\cdot\Delta L,\qquad\text{ligne repo}=s_b\cdot\Delta b,\qquad\text{ligne forward}=s_b\cdot\Delta b^F,
\]
et le \emph{total} est la somme des lignes dividende et repo.
\end{definition}

\begin{definition}[Exposition linéaire en forward]\label{def-fwdlinear}
Une position est linéaire en forward si $V(\theta)=\Phi\bigl(F(\theta)\bigr)$ pour une fonction $\Phi$ de classe $C^2$ :
sa valeur ne dépend du repo et des dividendes qu'à travers les forwards (options européennes avec une nappe de
volatilité qui n'en dépend pas, forwards, futures).
\end{definition}

\subsection{La jacobienne du repo implicite}
\begin{lemma}\label{lem-jacobian}
$J:=\partial_L\beta(F,L)=-(\partial_bF)^{-1}\partial_LF$. La matrice $\partial_bF$ est triangulaire inférieure avec des
termes diagonaux $\partial F_k/\partial b_k\le-T_kF_k<0$, donc $J$ existe. Au premier ordre,
\begin{equation}\label{eq-jacobian}
J_{ky}\approx-\frac{A_{ky}}{T_k\,(S-D_k)}.
\end{equation}
\end{lemma}
\begin{proof}
On différentie $F(\beta(F,L),L)=F$ par rapport à $L$ : $\partial_bF\,J+\partial_LF=0$. Avec un repo forward constant
entre piliers, $F_k$ ne dépend que de $b_1,\dots,b_k$, donc $\partial_bF$ est triangulaire inférieure ; sa diagonale vaut
$T_k\,\Phi_k'(B_k)\le-T_kF_k$ d'après le \cref{lem-monotone}. Pour le premier ordre, on fixe les forwards dans
\cref{eq-additivity} : $\delta B_k=-\delta D_k/(S-D_k)=-A_{ky}\delta L_y/(S-D_k)$, et on divise par $T_k$.
\end{proof}

\subsection{L'identité de contrôle}
\begin{theorem}[Total indépendant des dividendes]\label{thm-control}
Pour toute position,
\begin{equation}\label{eq-control}
\text{ligne dividende}+\text{ligne repo}=\text{ligne forward}+g\cdot\Delta L+O(\|\Delta\|^2),\qquad g=s_D+s_b\,J,
\end{equation}
avec $J$ évaluée en $(F^\T,L^\FO)$ et $\|\Delta\|=\|\Delta L\|+\|\Delta b^F\|$. Si la position est linéaire en forward,
alors $g=O(\|\Delta\|)$ : le total égale la ligne forward au second ordre près et ne dépend donc pas du panel
dividendes au premier ordre.
\end{theorem}
\begin{proof}
Un développement au premier ordre de $\beta$ en $L$ donne $\beta(F^\T,L^\T)=\beta(F^\T,L^\FO)+J\Delta L+O(\|\Delta L\|^2)$,
soit $\Delta b=\Delta b^F+J\Delta L+O(\|\Delta L\|^2)$. D'où
$s_D\cdot\Delta L+s_b\cdot\Delta b=s_b\cdot\Delta b^F+(s_D+s_bJ)\cdot\Delta L+O(\|\Delta\|^2)$.
Pour $V=\Phi\circ F$, $s_b=\nabla\Phi\,\partial_bF(\theta^\FO)$ et $s_D=\nabla\Phi\,\partial_LF(\theta^\FO)$, donc
$g=\nabla\Phi\,\bigl[\partial_LF+\partial_bF\,J\bigr](\theta^\FO)$. En $\theta'=(\beta(F^\T,L^\FO),L^\FO)$, le crochet est nul
d'après le \cref{lem-jacobian} ; comme $\theta'-\theta^\FO=(\Delta b^F,0)$ et que les dérivées sont lipschitziennes, le
crochet en $\theta^\FO$ est un $O(\|\Delta b^F\|)$, donc $g\cdot\Delta L=O(\|\Delta\|^2)$.
\end{proof}

\begin{theorem}[Exactitude de la réévaluation complète]\label{thm-exact}
Si la position est linéaire en forward et que toutes ses échéances sont des piliers Totem, le contrôle en
réévaluation complète $V(\theta^\T)-V(\theta^\FO)=\Phi(F^\T)-\Phi\bigl(F(\theta^\FO)\bigr)$ ne dépend pas du tout des
dividendes Totem $L^\T$.
\end{theorem}
\begin{proof}
Par construction du repo implicite, $F(\theta^\T)=F\bigl(\beta(F^\T,L^\T),L^\T\bigr)=F^\T$ quel que soit $L^\T$.
\end{proof}

\begin{remark}[Échéances entre piliers]
Pour une échéance située entre deux piliers, le forward dépend du portage interpolé et des dividendes alloués entre
les deux : une dépendance résiduelle à $L^\T$ subsiste. Elle est de l'ordre des dividendes détachés entre l'échéance et
les piliers voisins, c'est-à-dire exactement les effets de date et de profil des \cref{prop-timing,prop-profile}, que
le détecteur surveille.
\end{remark}

\begin{proposition}[Terme du second ordre]\label{prop-second}
Pour un pilier, avec $F(b,D)=e^{R-bT}(S-D)$ et $V=\Phi\circ F$, le total au premier ordre vaut
\begin{equation}\label{eq-second}
\Phi'(F)\,\Delta F-\Phi'(F)\,F\Bigl[\tfrac12T^2\Delta b^2+\frac{T\,\Delta b\,\Delta D}{S-D}\Bigr]+O(\|\Delta\|^3),
\end{equation}
alors que le contrôle exact vaut $\Phi'(F)\Delta F+\tfrac12\Phi''(F)\Delta F^2+O(\|\Delta\|^3)$. La dépendance du total au
premier ordre vis-à-vis du panel dividendes est donc entièrement contenue dans le crochet.
\end{proposition}
\begin{proof}
On développe $F$ au second ordre : $F_b=-TF$, $F_D=-F/(S-D)$, $F_{bb}=T^2F$, $F_{bD}=TF/(S-D)$, $F_{DD}=0$, d'où
$\Delta F=F_b\Delta b+F_D\Delta D+\tfrac12T^2F\Delta b^2+TF\Delta b\,\Delta D/(S-D)+O(3)$. Le total au premier ordre vaut
$\Phi'(F_b\Delta b+F_D\Delta D)$ ; on substitue.
\end{proof}

\subsection{Ce que la répartition coûte à un portefeuille}
\begin{proposition}[Matérialité de la répartition]\label{prop-materiality}
Supposons que l'erreur de consensus dividendes $\Delta L$ ait une covariance $\Sigma_L$ (bruit de consensus), les forwards
Totem étant fixés. Au premier ordre, la variance du total vaut $g^\top\Sigma_Lg$, contre $s_D^\top\Sigma_Ls_D$ pour la
seule ligne dividende. La part de l'incertitude de la ligne dividende compensée par la ligne repo vaut
$1-\sqrt{g^\top\Sigma_Lg\,/\,s_D^\top\Sigma_Ls_D}$.
\end{proposition}
\begin{proof}
D'après le \cref{thm-control}, le total est la ligne forward, qui ne fait pas intervenir $L^\T$, plus $g\cdot\Delta L$.
\end{proof}

\begin{corollary}[Ratio de compensation]
Pour un portefeuille linéaire en forward, un mouvement du panel dividendes produit des lignes dividende et repo
importantes et opposées, avec un total faible. Un ratio $1-|\text{total}|/(|\text{ligne div.}|+|\text{ligne repo}|)$
proche de un signale donc un artefact de répartition, pas une erreur de marquage.
\end{corollary}

$g$ se calcule à partir des seules sensibilités FO et de la jacobienne du bootstrap ; aucune revalorisation n'est
nécessaire.

\subsection{Exemples chiffrés}
\begin{example}[Un call]\label{ex-call}
Un call à la monnaie à un an (spot $100$, taux $3\,\%$, volatilité $25\,\%$, dividende FO $3{,}5$, repo FO $60$\,pb),
avec un forward Totem inférieur de $20$\,pb à celui du FO et trois panels dividendes différents :
\begin{center}\input{generated/\lang/tab_cancel}\end{center}
Le repo implicite varie de $125$\,pb d'un panel à l'autre et chaque ligne d'environ $0{,}3$, mais le total ne bouge pas.
\end{example}

\begin{example}[Trois portefeuilles]\label{ex-books}
Le pricer du prototype valorise trois portefeuilles sur onze piliers de repo et six périodes de dividendes : A (calls,
puts et un forward), B (A plus un swap de dividendes 2028 vendu) et C (A plus une commission de prêt de titres à un
an). Les forwards Totem sont inférieurs de $20$\,pb à ceux du FO et le panel dividendes Totem déplace les périodes
2027 et 2028 de $-0{,}4$, $0$ ou $+0{,}5$ :
\begin{center}\input{generated/\lang/tab_books}\end{center}
Pour le portefeuille A, le total reste à \vBookATotal{} alors que chaque ligne varie de plus de six unités. Pour B et C,
le total bouge, exactement de la ligne non-forward $g\cdot\Delta L$ : le swap de dividendes est une véritable
exposition aux dividendes et la commission de prêt une véritable exposition au repo.
\end{example}

\begin{figure}[ht]
\centering\includegraphics[width=\linewidth]{figures/\lang/cancellation.pdf}
\caption{Ligne dividende, ligne repo et leur total quand le panel dividendes Totem se déplace (portefeuille A à
gauche, B à droite). Points : réévaluation complète. Pour une exposition linéaire en forward, le total est plat
(\cref{thm-control,thm-exact}) ; le swap de dividendes ajoute une véritable exposition aux dividendes dont la ligne est
exactement $g\cdot\Delta L$.}
\label{fig-cancellation}
\end{figure}

\begin{example}[Matérialité]\label{ex-materiality}
Avec une dispersion de consensus dividendes de $5\,\%$ de chaque période, dix contributeurs et une corrélation
$e^{-|\Delta y|/2}$ entre périodes, la \cref{prop-materiality} donne
\begin{center}\input{generated/\lang/tab_materiality}\end{center}
\vMatA{} de l'incertitude de la ligne dividende se compense pour le portefeuille A ; \vMatB{} et \vMatC{} pour B et C.
\end{example}

\subsection{Quand la répartition compte}
L'exposition non-forward $g$ est non nulle pour les swaps et futures de dividendes (dividende pur), pour l'emprunt de
titres et les jambes de financement des swaps de rendement total (repo pur), pour l'exercice américain (la frontière
d'exercice anticipé dépend séparément des dividendes et du repo) et pour les exotiques valorisés avec une dynamique de
dividendes discrets (le modèle de dividendes modifie la volatilité du spot). Pour ces portefeuilles, la répartition
doit être contrôlée directement, et la qualité de la courbe de repo implicite (\cref{sec-detection}) compte.

\begin{insight}
\begin{itemize}[leftmargin=12pt]
\item Contrôler la \emph{ligne forward} et la \emph{ligne non-forward} : la première n'a besoin que des forwards Totem
et des dividendes FO ; la seconde n'a besoin des dividendes Totem que là où $g\neq0$.
\item Des lignes dividende et repo importantes qui se compensent, avec un total faible, sont un artefact de répartition.
\item En valorisation prudente~\cite{eu2016}, si les incertitudes sont agrégées par paramètre de valorisation sans
compensation, un artefact de répartition est compté deux fois ; définir le paramètre comme le forward là où $g\approx0$
supprime ce double comptage.
\end{itemize}
\end{insight}

\section{Un modèle statistique de l'écart de portage Totem--FO}\label{sec-noise}

\subsection{Modèle d'observation}
\begin{definition}[Modèle d'observation]\label{def-obs}
Le détecteur observe l'écart de portage $y\in\R^K$ de \cref{eq-y} et le modélise par
\begin{equation}\label{eq-obs}
y=X\beta+\sum_{t\in\mathcal{S}}a_t\,t+\varepsilon,\qquad X=\bigl[\,T,\;T\ln T\,\bigr],\qquad\varepsilon\sim\N(0,\kappa\Sigma),
\end{equation}
où $X\beta$ est un écart de repo lisse $\delta b(T)=\beta_0+\beta_1\ln T$ écrit dans l'espace de portage (désaccord
légitime, y compris une véritable erreur de marquage FO), $\mathcal{S}$ un ensemble éventuellement vide de gabarits
d'anomalie (\cref{sec-detection}), $\Sigma$ la covariance du bruit de consensus construite ci-dessous à partir des
statistiques Totem, et $\kappa>0$ une échelle unique calibrée sur un historique sain.
\end{definition}
$\beta_0$ est l'écart de repo à un an ($\ln1=0$) ; $\beta_1$ sa pente en $\ln T$. Sous l'hypothèse nulle,
$\mathcal{S}=\emptyset$.

\subsection{Covariance des moyennes de consensus}
D'après le \cref{lem-additivity} et la \cref{prop-gap}, le bruit de consensus en portage vaut
$\varepsilon_k=-e^F_k-e^D_k$, où $e^F_k$ est l'erreur du consensus du panel options $\ln\overline{F}_k$ et
$e^D_k=\delta\DB(T_k)/(S-D_k)$ celle du consensus dividendes alloué. Les deux panels étant indépendants
(\cref{ass-panels} avec $m=0$, le cas prudent),
\begin{equation}\label{eq-sigma}
\Sigma=\Sigma^F+\Sigma^D.
\end{equation}
Les deux termes sont des covariances de moyennes de panels emboîtés, donc la \cref{prop-nested} s'applique : chacun est
une covariance au niveau des contributeurs divisée terme à terme par $\max(n_k,n_l)$. Il reste à modéliser la
covariance au niveau des contributeurs à partir des dispersions publiées par Totem.

\subsection{Le panel options}
L'erreur de log-forward d'un contributeur en $T$ vaut $-T\,\delta b_j-\delta D_j(T)/(S-D)$. Les vues de repo sont lisses
(les desks marquent quelques piliers et interpolent) : la partie repo est essentiellement un niveau persistant, de
covariance $\sigma_b^2T_kT_l$. La partie dividendes est cumulative : un facteur persistant (un contributeur optimiste sur
les dividendes l'est pour toutes les années) plus des incréments qui arrivent avec chaque dividende.

\begin{lemma}[Covariance cumulative]\label{lem-cumcov}
Pour $s\in\R^K_{\ge0}$ croissant et $w\in[0,1]$, la matrice $C_{kl}=w\,s_ks_l+(1-w)\min(s_k,s_l)^2$ est semi-définie
positive : c'est la covariance de $Z_k=\sqrt{w}\,s_k\,\xi+\sqrt{1-w}\,W(s_k^2)$, avec $\xi\sim\N(0,1)$ indépendant d'un
mouvement brownien standard $W$. Pour $w=0$, la corrélation entre $k<l$ vaut $s_k/s_l$ (incréments indépendants).
\end{lemma}
\begin{proof}
$\Cov(Z_k,Z_l)=w\,s_ks_l+(1-w)\Cov\bigl(W(s_k^2),W(s_l^2)\bigr)=w\,s_ks_l+(1-w)\min(s_k^2,s_l^2)$, et une matrice de
covariance est semi-définie positive. Avec $w=0$, $\Cov(Z_k,Z_l)/\sqrt{\Var Z_k\Var Z_l}=s_k^2/(s_ks_l)=s_k/s_l$ pour $k<l$.
\end{proof}

\begin{proposition}[Covariance du panel options]\label{prop-sigmaF}
Soit $s_k=\sigma_{F,k}/F_k$ la dispersion publiée des forwards des contributeurs, en logarithme, et $\sigma_b$ la
dispersion des vues de repo. En décomposant $s_k^2=\sigma_b^2T_k^2+s_k'^2$,
\begin{equation}\label{eq-sigmaF}
\Sigma^F_{kl}=\frac{\sigma_b^2\,T_kT_l+w\,s'_ks'_l+(1-w)\min(s'_k,s'_l)^2}{\max(n_{A,k},n_{A,l})}.
\end{equation}
\end{proposition}
\begin{proof}
L'indépendance des vues de repo et de dividendes donne la somme d'un terme de rang un et de la covariance cumulative du
\cref{lem-cumcov} pour la partie dividendes ; la \cref{prop-nested} donne la division par le plus grand panel.
\end{proof}

\begin{remark}[Estimation de $\sigma_b$]
Aux échéances les plus courtes, l'incertitude sur les dividendes est négligeable, donc $s_k\approx\sigma_bT_k$. Le
prototype prend la médiane de $s_k/T_k$ sur les trois premiers piliers. Le minimum serait biaisé vers le bas (minimum de
trois estimations bruitées) : il sous-estimerait l'incertitude sur le niveau de repo et ferait déclencher trop souvent le
test d'erreur de marquage FO de la \cref{sec-detection}. Ce défaut a été constaté sur les titres difficiles à emprunter,
puis corrigé.
\end{remark}

\subsection{Le panel dividendes}
\begin{proposition}[Covariance du panel dividendes]\label{prop-sigmaD}
Supposons que les erreurs de consensus par période aient la covariance
$(\Sigma_L)_{yz}=\rho_{yz}\,\sigma_{D,y}\sigma_{D,z}/\max(n_{B,y},n_{B,z})$, avec $\rho_{yz}=e^{-|y-z|/\ell}$. Alors, au premier
ordre,
\begin{equation}\label{eq-sigmaD}
\Sigma^D=J_A\,\Sigma_L\,J_A^\top,\qquad (J_A)_{ky}=\frac{A_{ky}}{S-D_k}.
\end{equation}
\end{proposition}
\begin{proof}
$e^D_k=\delta D_k/(S-D_k)=\sum_yA_{ky}\delta L_y/(S-D_k)$ d'après le \cref{lem-allocation} : une application linéaire des
erreurs par période.
\end{proof}

\subsection{Les pics de dispersion sont des symptômes, pas du bruit}
\begin{proposition}[Masquage]\label{prop-masking}
Dans le désaccord de date de la \cref{prop-timing}, soit $s_k$ la dispersion hors désaccord et $\Delta=\tilde d_i/(S-D_k)$.
La statistique de test du déplacement en $T_k$, calculée avec la dispersion publiée (gonflée), vaut
\begin{equation}\label{eq-masking}
z=\frac{p\,\Delta}{\sqrt{s_k^2+p(1-p)\Delta^2}/\sqrt{n_{A,k}}},
\end{equation}
qui reste bornée par $\sqrt{n_{A,k}\,p/(1-p)}$ quelle que soit la taille du dividende, alors qu'avec la dispersion hors
désaccord $z$ croît linéairement en $\Delta$.
\end{proposition}
\begin{proof}
La \cref{prop-timing}(i) donne le déplacement moyen et (iii) la variance gonflée. Quand $\Delta\to\infty$, le rapport
tend vers $\sqrt{n_{A,k}}\,p/\sqrt{p(1-p)}$.
\end{proof}

\begin{example}
Pour l'exemple de la galerie ($p=0{,}37$, $n_A=16$), la borne vaut $\sqrt{16\times0{,}37/0{,}63}\approx3{,}1$ : à peine
le seuil critique de la famille de tests, si bien que le désaccord passerait souvent inaperçu. Le prototype plafonne donc
un pic isolé de dispersion à $1{,}5$ fois la plus grande des deux dispersions voisines avant de construire $\Sigma$, et
publie le rapport, qui est lui-même un indice de désaccord entre contributeurs, en général sur une date de détachement
(ici $9{,}7$).
\end{example}

\begin{insight}
$\Sigma$ n'utilise que ce que Totem publie (dispersions et nombres de contributeurs) plus trois réglages structurels :
le poids du facteur $w$, la longueur de corrélation $\ell$ entre périodes de dividendes et la dispersion du repo lue sur
le court terme. L'échelle unique $\kappa$ absorbe un biais commun de ces réglages ; c'est leur \emph{forme} qui rend le
détecteur puissant, car elle lui indique quels écarts sont bon marché (lisses, longs) et lesquels sont coûteux (courts,
locaux).
\end{insight}

\section{Détection : filtres adaptés et test du rapport de vraisemblance}\label{sec-detection}

\subsection{Gabarits}
Chaque mécanisme de la \cref{sec-anatomy} laisse une forme connue dans l'espace de portage. Le détecteur teste un
\emph{gabarit} $t\in\R^K$ par mécanisme et par position :
\begin{center}\small
\begin{tabular}{lp{3.3cm}p{4.3cm}p{4.6cm}}
\toprule
Gabarit & Vecteur & Mécanismes & Amplitude $a$ en unités naturelles\\
\midrule
niveau & $\one$ & décalage du spot ; dividende avant $T_1$ & erreur de spot $\approx-a$ ; dividende proche $\approx-a(S-D_1)$\\
marche $k$ & $\one\{T\ge T_k\}$ & montant de dividende dans $(T_{k-1},T_k]$ ; changement de panel en $T_k$ & erreur de dividende $\approx-a(S-D_k)$\\
pic $k$ & $e_k$ & forward aberrant ; date de détachement de part et d'autre de $T_k$ & $\delta\ln F_k\approx-a$ ; part du dividende $\approx-a(S-D_k)/\tilde d_i$\\
saison & $\sum_i\tilde d_iK_h(T_k-t_i)/(S-D_k)$ & décalage du profil de détachement & décalage $\approx365\,a$ jours\\
\bottomrule
\end{tabular}
\end{center}
$K_h$ est un noyau gaussien ($h=20$ jours). Les unités naturelles découlent des
\cref{prop-spot,prop-bucket,prop-timing,prop-profile}. Un pic au dernier pilier est le même vecteur qu'une marche à ce
pilier et il est présenté comme l'un ou l'autre. Avec onze piliers, on obtient \vNTemplates{} gabarits.

\subsection{Blanchiment, projection et filtre adapté}
Soit $\kappa\Sigma=LL^\top$ (Cholesky), $\tilde y=L^{-1}y$, $\tilde X=L^{-1}X$, et $P=I-QQ^\top$ avec $Q$ une base
orthonormée des colonnes de $\tilde X$ : $P$ retire le modèle nul dans l'espace blanchi. On note $r=P\tilde y$ et
$\tb=PL^{-1}t$.

\begin{theorem}[Test du rapport de vraisemblance généralisé pour un gabarit]\label{thm-glrt}
Dans le modèle $y=X\beta+a\,t+\varepsilon$, $\varepsilon\sim\N(0,\kappa\Sigma)$, avec $\beta$ inconnu, la statistique du
rapport de vraisemblance généralisé de $H_0:a=0$ contre $H_1:a\ne0$ est $z_t^2$, avec
\begin{equation}\label{eq-z}
z_t=\frac{\tb^\top r}{\|\tb\|},\qquad z_t\sim\N\bigl(a\|\tb\|,\,1\bigr).
\end{equation}
L'amplitude du maximum de vraisemblance est $\hat a_t=\tb^\top r/\|\tb\|^2$, d'écart-type $1/\|\tb\|$. Parmi les tests
invariants par le sous-espace du modèle nul et par le signe de $a$, le test qui rejette pour les grandes valeurs de
$|z_t|$ est uniformément le plus puissant.
\end{theorem}
\begin{proof}
Après blanchiment, $\tilde y=\tilde X\beta+a\,L^{-1}t+\tilde\varepsilon$ avec $\tilde\varepsilon\sim\N(0,I)$. La
log-vraisemblance maximisée vaut $-\frac12$ fois la somme des carrés des résidus, soit $\|P\tilde y\|^2$ sous $H_0$ et le
résidu de la régression sur $[\tilde X,L^{-1}t]$ sous $H_1$. D'après le théorème de Frisch--Waugh--Lovell~\cite{lovell1963},
ce dernier vaut $\|r\|^2-(\tb^\top r)^2/\|\tb\|^2$, donc $2\ln\Lambda=(\tb^\top r)^2/\|\tb\|^2=z_t^2$, et $\hat a_t$ est le
coefficient de la régression partielle. Comme $P\tilde X=0$ et $\tb^\top P=\tb^\top$,
$\tb^\top r=a\,\tb^\top L^{-1}t+\tb^\top\tilde\varepsilon=a\|\tb\|^2+\tb^\top\tilde\varepsilon$, et
$\tb^\top\tilde\varepsilon/\|\tb\|\sim\N(0,1)$. L'optimalité est le résultat des détecteurs de sous-espace adapté de
Scharf et Friedlander~\cite{scharf1994} ; voir aussi~\cite{kay1998}.
\end{proof}

\begin{corollary}[Puissance et rapport signal sur bruit]\label{cor-power}
Au seuil $c$, la probabilité de détection vaut $\Phi(\lambda-c)+\Phi(-\lambda-c)$, où $\lambda=|a|\,\|\tb\|$ est le rapport
signal sur bruit de l'anomalie. Avec le seuil familial typique du banc d'essai, $c\approx\vZcrit$, une anomalie de
$\lambda=4$ est détectée avec une probabilité de \vPowThFour{} et une anomalie de $\lambda=6$ avec \vPowThSix.
\end{corollary}
\begin{proof}
$|z_t|>c$ avec $z_t\sim\N(\lambda,1)$ (au signe près).
\end{proof}

\begin{proposition}[Test global]\label{prop-q}
Sous $H_0$, $Q=\|r\|^2\sim\chi^2_{K-p}$, où $p=2$ est le nombre de colonnes du modèle nul.
\end{proposition}
\begin{proof}
$r=P\tilde\varepsilon$, où $P$ est une projection orthogonale de rang $K-p$.
\end{proof}
$Q$ détecte une rugosité qu'aucun gabarit n'explique.

\subsection{Calibration familiale}
\begin{proposition}[Valeur critique familiale]\label{prop-family}
Sous $H_0$, le vecteur $z=(z_t)_t$ est gaussien centré de corrélation $C_{tu}=\tb^\top\tilde u/(\|\tb\|\|\tilde u\|)$, qui ne
dépend ni de $\beta$ ni de $\kappa$. La valeur critique $c_\alpha$ telle que $P(\max_t|z_t|>c_\alpha)=\alpha/2$ se calcule
exactement en simulant $\N(0,C)$ ; elle est au plus égale à la valeur de Šidák
$\Phi^{-1}\bigl(1-\tfrac12[1-(1-\alpha/2)^{1/m}]\bigr)$ pour $m$ gabarits. La règle « alerte si $\max_t|z_t|>c_\alpha$ ou si
$Q$ dépasse son quantile $1-\alpha/2$ » a un taux de fausses alertes d'au plus $\alpha$.
\end{proposition}
\begin{proof}
$z=D^{-1}\tilde T^\top P\tilde\varepsilon$, où $\tilde T$ est la matrice des $\tb$ et $D$ celle de leurs normes, d'où la
corrélation. L'inégalité de Šidák~\cite{sidak1967} donne $P(\max_t|z_t|\le c)\ge\prod_tP(|z_t|\le c)$ pour toute
corrélation : le seuil de Šidák est donc conservateur. Les deux tests se partagent le budget $\alpha$ par l'inégalité de
la réunion.
\end{proof}

\begin{proposition}[Calibration exacte de l'échelle de variance]\label{prop-kappa}
Remplacer $\kappa=1$ par $\kappa$ divise chaque $z_t$ par $\sqrt\kappa$ et $Q$ par $\kappa$, et laisse $c_\alpha$ inchangé.
Pour un ensemble de sous-jacents sains, l'indicatrice d'alerte est donc décroissante en $\kappa$, le taux d'alerte
empirique est une fonction en escalier décroissante de $\kappa$, et le plus petit $\kappa$ dont le taux d'alerte ne dépasse
pas la cible s'obtient exactement par dichotomie sur $\ln\kappa$.
\end{proposition}
\begin{proof}
Blanchir avec $\sqrt\kappa L$ multiplie $\tilde y$, $\tilde X$ et chaque $\tb$ par $\kappa^{-1/2}$ ; $P$ est inchangé, donc $r$ et
$\tb$ sont multipliés par $\kappa^{-1/2}$, $z_t$ par $\kappa^{-1/2}$ et $Q$ par $\kappa^{-1}$ ; $C$ est une corrélation et ne bouge
pas. L'indicatrice $\{\max|z|/\sqrt\kappa>c\}\cup\{Q/\kappa>q\}$ de chaque sous-jacent se réduit quand $\kappa$ augmente.
\end{proof}

\subsection{Attribution, réparation et test d'erreur de marquage FO}
Après une alerte, les gabarits sont retenus de façon gloutonne (poursuite adaptative~\cite{mallat1993}) : à chaque
étape, on ajoute le gabarit de plus grand $|z|$ après projection du modèle nul et des gabarits déjà retenus, tant qu'il
dépasse $c_\alpha$. Les gabarits retenus sont ensuite réajustés conjointement avec le modèle nul par moindres carrés
généralisés, chaque amplitude est traduite en unités naturelles, et le portage \emph{réparé} vaut
$y-\sum_{t\in\mathcal S}\hat a_t\,t$. Les gabarits contenus dans le sous-espace du modèle nul sont indétectables par
construction : un écart de repo lisse est un désaccord légitime, publié mais jamais signalé comme anomalie.

Le niveau de repo $\hat\beta_0$ (Totem moins FO à un an) et son écart-type proviennent du même ajustement. Le repo FO est
signalé comme erreur de marquage possible lorsque $|\hat\beta_0|$ dépasse à la fois $10$\,pb et
$t_{\nu,1-\alpha/2}\cdot s\cdot\operatorname{se}(\hat\beta_0)$ :
\begin{itemize}
\item $t_\nu$ est un quantile de Student à $\nu=n_A-1$ degrés de liberté sur le court terme : l'écart-type repose sur des
dispersions estimées à partir d'une poignée de contributeurs, et une estimation gaussienne divisée par une échelle
indépendante de type $\chi^2_\nu$ suit une loi de Student ;
\item $s\ge1$ est une échelle calibrée sur un historique sain exactement comme $\kappa$, de sorte que les sous-jacents
sains déclenchent au plus $\alpha$ alertes d'erreur de marquage FO (\vLevelScaleEuStock, \vLevelScaleUsStock,
\vLevelScaleHtb{} et \vLevelScaleIndex{} pour les quatre catégories du banc d'essai).
\end{itemize}

\begin{example}[Une détection à la main]\label{ex-toy}
Quatre piliers $T=(0{,}25\,;0{,}5\,;1\,;2)$, un bruit indépendant d'écarts-types $(1\,;2\,;4\,;8)$\,pb de portage, un modèle
nul réduit au niveau de repo ($X=T$), un niveau vrai de $30$\,pb, un pic de $12$\,pb en $T_2$ et un bruit
$(0{,}4\,;-1{,}1\,;2{,}6\,;-3{,}0)$\,pb, d'où $y=(7{,}9\,;25{,}9\,;32{,}6\,;57{,}0)$\,pb. Le niveau ajusté vaut
$\hat\beta_0=\vToyBeta$\,pb et $Q=\vToyQ$ (le quantile à $99\,\%$ de $\chi^2_3$ vaut $11{,}3$). Les \vToyM{} gabarits donnent
\begin{center}\input{generated/\lang/tab_toy}\end{center}
La valeur critique familiale à $\alpha=2\,\%$ vaut $c=\vToyC$ (Šidák : \vToyCsidak). Le pic en $T_2$ l'emporte
($z=4{,}52$, amplitude $10{,}5\pm2{,}3$\,pb pour $12$ injectés) ; la marche à partir de $T_2$ obtient aussi un score élevé car
les deux gabarits sont corrélés. La poursuite adaptative retient le pic, après quoi tous les $|z|$ restants sont inférieurs
à $0{,}8$ et $Q$ tombe à $0{,}6$ : une seule explication suffit.
\end{example}

\subsection{Algorithmes}
\begin{algorithm}[ht]
\caption{Bootstrap du repo implicite}\label{alg-bootstrap}
\begin{algorithmic}[1]
\Require spot $S$, taux cumulé $R(\cdot)$, dividendes $(t_i,d_i)$, piliers $T_1<\dots<T_K$, forwards $F_1,\dots,F_K$
\For{$k=1,\dots,K$}
  \State construire $\Phi_k$ de \cref{eq-phik} à partir de $B_1,\dots,B_{k-1}$
  \State $x\gets R_k-\ln\bigl(F_k/(S-\hat D_k)\bigr)$ \Comment{dividendes de la période au repo précédent}
  \Repeat
    \State $g\gets\ln\Phi_k(x)-\ln F_k$ ; mettre à jour l'encadrement $[\ell,u]$ selon le signe de $g$
    \State $x\gets x-g\,\Phi_k(x)/\Phi_k'(x)$ ; si $x\notin(\ell,u)$, $x\gets(\ell+u)/2$
  \Until{le pas est inférieur à la tolérance}
  \State $B_k\gets x$
\EndFor
\State \Return $b_k=B_k/T_k$
\end{algorithmic}
\end{algorithm}

\begin{algorithm}[ht]
\caption{Détection structurelle}\label{alg-detect}
\begin{algorithmic}[1]
\Require écart de portage $y$, covariance $\Sigma$, échelle $\kappa$, modèle nul $X$, gabarits $\mathcal T$, niveau $\alpha$
\State $L\gets\operatorname{chol}(\kappa\Sigma)$ ; $P\gets$ projecteur qui retire $L^{-1}X$ ; $r\gets PL^{-1}y$ ; $Q\gets\|r\|^2$
\State pour chaque $t\in\mathcal T$ : $\tb\gets PL^{-1}t$, $z_t\gets\tb^\top r/\|\tb\|$
\State $c_\alpha\gets$ quantile $1-\alpha/2$ de $\max_t|Z_t|$, $Z\sim\N(0,C)$ \Comment{\cref{prop-family}}
\State alerte $\gets\max_t|z_t|>c_\alpha$ \textbf{ou} $Q>\chi^2_{K-2}(1-\alpha/2)$
\If{alerte}
  \State $\mathcal S\gets\emptyset$ ; \textbf{répéter} ajouter le gabarit de plus grand $|z|$ après projection de $[X,\mathcal S]$
  \Statex \hspace{\algorithmicindent} \textbf{tant que} ce $|z|$ dépasse $c_\alpha$
\EndIf
\State ajustement MCG conjoint de $y$ sur $[X,\mathcal S]$ ; amplitudes en unités naturelles ; portage réparé $y-\sum\hat a_tt$
\State alerte d'erreur de marquage FO si $|\hat\beta_0|>\max\bigl(10\text{\,pb},\;t_{\nu}\,s\operatorname{se}(\hat\beta_0)\bigr)$
\end{algorithmic}
\end{algorithm}

\begin{algorithm}[ht]
\caption{Calibration de l'échelle de variance}\label{alg-kappa}
\begin{algorithmic}[1]
\Require sous-jacents sains $i=1,\dots,N$ avec $(\max_t|z_{t,i}|,\,c_i,\,Q_i)$ calculés pour $\kappa=1$, cible $\alpha$
\State $\ell\gets\ln10^{-3}$, $u\gets\ln10^{3}$
\For{$60$ itérations}
  \State $m\gets(\ell+u)/2$ ; taux $\gets\frac1N\sum_i\one\bigl\{\max|z_i|e^{-m/2}>c_i\ \text{ou}\ Q_ie^{-m}>q_{1-\alpha/2}\bigr\}$
  \State \textbf{si} taux $\le\alpha$ \textbf{alors} $u\gets m$ \textbf{sinon} $\ell\gets m$
\EndFor
\State \Return $\kappa=e^{u}$
\end{algorithmic}
\end{algorithm}

\subsection{Exemples sur des sous-jacents synthétiques}
La \cref{fig-decomposition} montre le mécanisme dans l'espace de portage pour une erreur sur une période de dividendes,
et la \cref{fig-gallery} six détections dans l'espace de repo, avec les explications du détecteur au
\cref{tab-gallery}.

\begin{figure}[ht]
\centering\includegraphics[width=\linewidth]{figures/\lang/decomposition.pdf}
\caption{Vue en espace de portage de l'erreur sur une période de dividendes Totem de la \cref{fig-gallery}(b) : écart de
portage observé, ajustement nul avec sa bande de $\pm2$ écarts-types, et ajustement nul plus le gabarit de marche retenu.}
\label{fig-decomposition}
\end{figure}

\begin{figure}[p]
\centering\includegraphics[width=\linewidth]{figures/\lang/gallery.pdf}
\caption{Six actions européennes synthétiques, une par scénario, avec les réglages du détecteur du banc d'essai : repo
FO, repo implicite Totem, repo implicite après retrait de l'anomalie détectée, et bande de bruit par pilier. En (f), la
courbe Totem est saine et le repo FO est décalé : aucune anomalie n'est signalée et l'erreur de marquage FO est
publiée ; la bande est par pilier alors que le test d'erreur de marquage regroupe tous les piliers.}
\label{fig-gallery}
\end{figure}

\begin{table}[ht]
\centering
\caption{Première explication fournie par le détecteur pour chaque panneau de la \cref{fig-gallery}.}\label{tab-gallery}
\input{generated/\lang/tab_gallery}
\end{table}

\subsection{Les petits réseaux de neurones sont des filtres adaptés}
\begin{proposition}\label{prop-nn}
Supposons le modèle nul projeté et le bruit blanchi, de sorte que $r\sim\N(0,P)$ sous $H_0$ et $r\sim\N(a\tb,P)$ sous $H_1$,
avec un gabarit $t$ et une amplitude $a$ connus. Le logarithme du rapport des probabilités a posteriori vaut
\begin{equation}\label{eq-nn}
\ln\frac{P(H_1\mid r)}{P(H_0\mid r)}=a\,\tb^\top r-\tfrac12a^2\|\tb\|^2+\ln\frac{\pi_1}{\pi_0},
\end{equation}
une fonction affine du filtre adapté $\tb^\top r$. Un classifieur entraîné par entropie croisée sur suffisamment de données
converge vers cette probabilité a posteriori : dans ce cadre, un « petit réseau par type d'anomalie » ne peut au mieux
qu'apprendre le filtre adapté ; avec une amplitude et un signe inconnus, on retrouve la statistique du rapport de
vraisemblance $|z_t|$ du \cref{thm-glrt}.
\end{proposition}
\begin{proof}
Les deux densités gaussiennes ont la même covariance $P$, donc leur rapport logarithmique est affine en $r$ : les termes
quadratiques se compensent. L'entropie croisée est minimisée par la vraie probabilité a posteriori. Maximiser
\cref{eq-nn} en $a$ donne $(\tb^\top r)^2/(2\|\tb\|^2)$.
\end{proof}

\begin{insight}
Le détecteur est le détecteur optimal pour les mécanismes que l'on sait nommer, écrit plutôt qu'appris : son taux de
fausses alertes est fixé exactement, sa puissance est connue en forme fermée (\cref{cor-power}), et chaque alerte est
accompagnée d'une amplitude en unités naturelles. Les modèles appris trouvent leur place sur ce que les gabarits ne
couvrent pas (\cref{sec-alternatives}).
\end{insight}

\section{Banc d'essai sur un marché Totem synthétique}\label{sec-benchmark}

\subsection{Conception}
Le marché synthétique est construit à partir des contributeurs, pour que chaque mécanisme de la \cref{sec-anatomy}
apparaisse comme dans Totem et non comme une forme ajoutée à une courbe.

\paragraph{Catégories.} Quatre catégories, onze échéances Totem d'un mois à cinq ans et six périodes de dividendes (la
fin de 2026, puis 2027 à 2031) ; date de valorisation au 14 octobre 2026.
\begin{center}\input{generated/\lang/tab_categories}\end{center}

\paragraph{Contributeurs.} Les panels options et dividendes sont tirés indépendamment. Chaque contributeur a des vues
cohérentes (\cref{ass-consistent}) : des dividendes avec un facteur persistant plus un bruit propre à chaque dividende,
dont la dispersion croît avec le temps, et une courbe de repo lisse (niveau plus inclinaison). Les contributeurs options
valorisent calls et puts à cinq strikes avec leur propre forward, une nappe de volatilité avec smile, un bruit de
volatilité et un bruit de prix. Les panels se réduisent avec la maturité et sont emboîtés.

\paragraph{Consensus et chaîne de traitement.} Le consensus est la moyenne après un filtre médiane/MAD ; dispersions et
nombres de contributeurs sont publiés comme le ferait Totem. La chaîne du contrôleur s'exécute ensuite sans
modification : forwards par parité, allocation sur le profil de détachement FO, bootstrap, covariance, détection.

\paragraph{Marquages FO.} Dividendes vrais avec un bruit de $1\,\%$ et repo vrai avec une erreur de niveau de $3$\,pb (un
écart-type).

\paragraph{Scénarios.} Chaque sous-jacent de test est simulé sous chaque scénario avec la même graine : les différences
entre scénarios sont appariées.
\begin{center}\small
\begin{tabular}{p{4.0cm}p{7.4cm}p{3.2cm}}
\toprule
Scénario & Injection & Amplitude\\
\midrule
Décalage du spot & les contributeurs options valorisent avec un autre spot & $\pm0{,}03\,\%$ à $\pm0{,}15\,\%$\\
Erreur de période & une période (2027--2029) erronée dans le consensus dividendes & $\pm0{,}15$ à $\pm0{,}5$ pour $100$ de spot\\
Désaccord de date & $35$ à $70\,\%$ du panel options placent le dividende le plus proche d'une échéance de l'autre côté & un dividende\\
Profil périmé & profil d'allocation décalé pour qu'un dividende franchisse une échéance & $3$ à $20$ jours au-delà\\
Forward aberrant & calls de consensus décalés à une échéance & $3$ à $6$ écarts-types du bruit de portage\\
Sous-panel long biaisé & les contributeurs restant après la plus forte sortie sont biaisés sur les dividendes suivants & $\pm8\,\%$ à $\pm20\,\%$\\
Erreur de marquage FO & repo FO décalé parallèlement, Totem intact & $\pm15$ à $\pm50$\,pb\\
Biais brut/net & toutes les périodes Totem multipliées & $\pm8\,\%$ à $\pm18\,\%$\\
\bottomrule
\end{tabular}
\end{center}

\paragraph{Détecteurs.} (1) Écart naïf : plus grand $|\Delta b|$ sur les piliers. (2) Forme naïve : plus grand écart du
repo implicite Totem à l'interpolation de ses deux voisins. (3) B-spline plus mélange gaussien sur la courbe Totem :
coefficients d'une B-spline cubique ajustée par moindres carrés (six coefficients, nœuds en $\ln T$), regroupés par un
mélange gaussien par catégorie dont le nombre de composantes est choisi par BIC ; le score est l'opposé de la
log-vraisemblance. (4) La même chose sur l'écart Totem--FO. (5) Le détecteur structurel de la \cref{sec-detection}.

\paragraph{Protocole.} \vNTrain{} sous-jacents sains par catégorie calibrent chaque détecteur à \vFPR{} de fausses
alertes (un quantile du score pour les références, $\kappa$ et l'échelle de niveau pour le détecteur structurel).
\vNTestPerCat{} nouveaux sous-jacents par catégorie sont ensuite simulés sous chacun des \vNScen{} scénarios, soit
\vNTest{} sous-jacents de test. Le plan étant apparié, le rapport signal sur bruit d'une injection se calcule exactement,
$\lambda=\|PL^{-1}(y^{\text{scénario}}-y^{\text{sain}})\|$ : c'est la grandeur du \cref{cor-power}.

\subsection{Résultats}
\begin{table}[ht]
\centering
\caption{Part des sous-jacents signalés, chaque détecteur étant calibré à \vFPR{} de fausses alertes sur l'historique
sain. La ligne « sain » est le taux de fausses alertes réalisé sur de nouveaux sous-jacents. Pour l'erreur de marquage FO
et le biais brut/net, une alerte n'est pas le résultat souhaité (voir le texte). Colonnes : écart naïf, forme naïve,
B-spline plus mélange gaussien sur la courbe et sur l'écart Totem--FO, détecteur structurel.}\label{tab-detection}
\input{generated/\lang/tab_detection}
\end{table}

\begin{figure}[ht]
\centering\includegraphics[width=\linewidth]{figures/\lang/detection.pdf}
\caption{Part des sous-jacents signalés par scénario et par détecteur (données du \cref{tab-detection}).}
\label{fig-detection}
\end{figure}

\begin{figure}[ht]
\centering\includegraphics[width=\linewidth]{figures/\lang/power.pdf}
\caption{Taux de détection en fonction du rapport signal sur bruit $\lambda$ de l'anomalie injectée, sur l'ensemble des six
anomalies côté Totem ; $n$ est le nombre de sous-jacents par classe. La courbe grise est la puissance théorique d'un
gabarit connu au seuil familial typique (\cref{cor-power}).}
\label{fig-power}
\end{figure}

\begin{table}[ht]
\centering
\caption{Aire sous la courbe ROC par rapport aux sous-jacents sains (qualité du classement, indépendante du seuil).}\label{tab-auc}
\input{generated/\lang/tab_auc}
\end{table}

\begin{table}[ht]
\centering
\caption{Détecteur structurel : part des sous-jacents signalés par scénario et par catégorie.}\label{tab-bycat}
\input{generated/\lang/tab_bycat}
\end{table}

\begin{figure}[ht]
\centering\includegraphics[width=\linewidth]{figures/\lang/attribution.pdf}
\caption{Mécanisme cité en premier par le détecteur structurel, par scénario. Niveau, marche, pic et saison sont les
familles de gabarits ; « rugosité » signifie que le test global s'est déclenché sans gabarit ; « pas d'alerte » signifie
que le sous-jacent n'a pas été signalé.}
\label{fig-attribution}
\end{figure}

\begin{table}[ht]
\centering
\caption{Détecteur structurel : détection, attribution (la première explication cite le mécanisme injecté et sa
position) et alertes d'erreur de marquage FO.}\label{tab-attribution}
\input{generated/\lang/tab_attribution}
\end{table}

\begin{table}[ht]
\centering
\caption{Taux de fausses alertes réalisé sur de nouveaux sous-jacents sains, par catégorie (cible \vFPR), et $\kappa$
calibré.}\label{tab-fpr}
\input{generated/\lang/tab_fpr}
\end{table}

\subsection{Lecture des résultats}
\begin{itemize}
\item \textbf{La détection moyenne ne départage pas les meilleures méthodes.} Sur les six anomalies côté Totem, le
détecteur structurel signale \vMeanStruct{} des sous-jacents, la B-spline plus mélange gaussien sur l'écart
\vMeanGmmDiff, sur la courbe \vMeanGmmCurve, le contrôle de forme naïf \vMeanShape{} et l'écart naïf \vMeanNaive.
\item \textbf{Les forwards aberrants} sont le cas où l'espace de portage et le modèle de bruit rapportent le plus :
\vFwdStruct{} détectés, contre \vFwdNaive{} pour l'écart naïf, qui ne peut pas voir une valeur aberrante de trois à six
écarts-types à une échéance longue sous le bruit du long terme, et \vFwdShape{} pour le contrôle de forme naïf.
\item \textbf{Les décalages de spot} (\vSpotStruct) et \textbf{les profils de détachement périmés} (\vSeasonStruct) sont
faciles pour toutes les méthodes ; le détecteur structurel ajoute l'explication (\cref{fig-attribution}).
\item \textbf{Les erreurs de période} (\vBucketStruct{} contre \vBucketNaive{} pour l'écart naïf) sont limitées par leur
taille : une erreur de $0{,}15$ à $0{,}5$ pour $100$ de spot est faible devant la dispersion des dividendes lointains, et la
\cref{fig-power} montre que les cas manqués ont un faible rapport signal sur bruit.
\item \textbf{Les désaccords de date de détachement} sont détectés dans \vTimingStruct{} des cas ; le contrôle de forme
naïf (\vTimingShape) est compétitif car ces pics sont grands dans l'espace de repo.
\item \textbf{Un sous-panel long biaisé} est presque non identifiable pour toutes les méthodes (\vPanelStruct{} pour le
détecteur structurel, \vPanelNaive{} pour l'écart naïf) : le biais s'accumule de façon lisse après la sortie de
contributeurs et se projette surtout sur une pente de repo. C'est le \cref{thm-nonid} à l'œuvre, pas un problème de
réglage.
\item \textbf{Les véritables erreurs de marquage FO} restent distinctes des courbes Totem aberrantes : \vFoStruct{}
déclenchent une anomalie (le taux sain) et \vFoMismark{} sont signalées comme erreurs de marquage FO possibles, contre
\vFoClean{} des sous-jacents sains. L'écart naïf et le mélange gaussien sur l'écart déclenchent une alerte de courbe
aberrante sur \vFoNaive{} et \vFoGmmDiff{} d'entre elles. Le test d'erreur de marquage FO trouve \vFoMismarkEuStock{} des
erreurs sur les actions européennes et \vFoMismarkIndex{} sur l'indice, mais seulement \vFoMismarkHtb{} sur les titres
difficiles à emprunter : les vues de repo des contributeurs y diffèrent d'environ $60$\,pb, si bien qu'un écart FO de
$15$ à $50$\,pb reste dans l'incertitude du consensus lui-même.
\item \textbf{Le biais brut/net} est signalé dans \vGrossStruct{} des cas, surtout sous forme de marches aux dates de
détachement pour les payeurs discrets ; pour le flux de dividendes d'un indice, il ressemble à une pente de repo et
ressort via $\beta$ (\cref{prop-gross}).
\item \textbf{Puissance.} Entre $4$ et $6$ de rapport signal sur bruit, le détecteur structurel trouve \vPowStructE{} des
anomalies, et au-delà de $6$ de \vPowStructF{} à \vPowStructH, conformément à la théorie du \cref{cor-power}. En dessous de
$3$, l'écart naïf et le mélange gaussien sur l'écart détectent davantage (\vPowGmmDiffC{} contre \vPowStructC{} entre $2$
et $3$) : ils réagissent aussi à la part d'une anomalie qui ressemble à un écart de repo lisse, que le détecteur
structurel laisse volontairement au test d'erreur de marquage FO. C'est cette même propriété qui leur fait déclencher des
alertes sur les véritables erreurs de marquage FO.
\item \textbf{Attribution.} Quand le détecteur structurel signale une anomalie Totem, sa première explication cite le
mécanisme injecté et sa position dans \vAttr{} des cas.
\end{itemize}

\subsection{Limites des éléments présentés}
\begin{caveat}
Le marché synthétique reproduit les mécanismes par construction ; il ne montre pas que les données Totem réelles ne
contiennent que ceux-là. Les données réelles comporteront aussi des contributeurs figés, des opérations sur titres et des
dividendes exceptionnels, de l'exercice américain dans les cotations d'actions et des erreurs de saisie. L'échelle $\kappa$
absorbe un niveau de bruit mal spécifié, mais pas une structure mal spécifiée, et les amplitudes, bien que plausibles, ne
sont calibrées sur aucun historique particulier. Le banc d'essai montre que la méthode fonctionne quand les mécanismes sont
présents et qu'elle est calibrée quand ils sont absents ; la validation sur douze à vingt-quatre mois d'instantanés réels
avec des incidents étiquetés est l'étape suivante.
\end{caveat}

\section{La place des B-splines, des mélanges gaussiens et des réseaux de neurones}\label{sec-alternatives}

\subsection{Ce que les coefficients B-spline ne voient pas}
\begin{proposition}[Levier d'un pic]\label{prop-leverage}
Soit $B\in\R^{K\times q}$ la matrice de conception B-spline aux piliers et $H=B(B^\top B)^{-1}B^\top$ la matrice chapeau.
L'ajustement par moindres carrés d'un pic unitaire $e_k$ a une énergie $\|He_k\|^2=h_{kk}$, et $\sum_kh_{kk}=q$. Une
fraction $1-h_{kk}$ de l'énergie du pic reste dans le résidu et est invisible pour toute variable calculée à partir des
coefficients.
\end{proposition}
\begin{proof}
$H$ est symétrique et idempotente, donc $\|He_k\|^2=e_k^\top H^\top He_k=e_k^\top He_k=h_{kk}$, et
$\sum_kh_{kk}=\operatorname{tr}H=\rank B=q$. Les coefficients $(B^\top B)^{-1}B^\top e_k$ déterminent $He_k$ et rien d'autre.
\end{proof}

\begin{example}
Pour les six B-splines cubiques (nœuds en $\ln T$) du banc d'essai et les onze piliers,
\begin{center}\input{generated/\lang/tab_leverage}\end{center}
Aux piliers intérieurs, entre neuf mois et quatre ans, les coefficients ne captent que \vLevMin{} à \vLevMax{} de
l'énergie d'un pic (en moyenne $6/11\approx\num{0.545}$ sur l'ensemble des piliers). La spline répartit en outre ce qu'elle
capte sur plusieurs coefficients. Une anomalie sur un seul pilier, le motif de courbe aberrante le plus fréquent, est donc
diluée avant même que le regroupement ne commence.
\end{example}

\subsection{Pourquoi un mélange sur l'écart signale les erreurs de marquage FO}
\begin{proposition}\label{prop-gmm-fo}
Soient les variables les coefficients spline $c=B^+\Delta b$ de l'écart de repo. Un décalage parallèle du repo FO de
$\delta$ les déplace de $-\delta\,B^+\one$, une direction fixe. Si le mélange est ajusté sur un historique sain, où FO et Totem
concordent au bruit près, sa densité décroît dans cette direction : son taux d'alerte sur les erreurs de marquage FO tend
vers un quand $|\delta|$ augmente, quelle que soit la qualité de la courbe Totem.
\end{proposition}
\begin{proof}
$c$ est linéaire en $\Delta b$, et $\Delta b$ se déplace de $-\delta\one$. Chaque composante gaussienne a une variance finie dans
la direction $B^+\one$, donc la densité du mélange, et avec elle le score de vraisemblance, tend vers zéro dans cette
direction.
\end{proof}
Le détecteur structurel n'a pas ce problème : $T\cdot\one$ appartient au sous-espace du modèle nul, donc un écart de niveau
est retiré de chaque test et publié via $\hat\beta_0$. Le banc d'essai confirme les deux énoncés : \vFoGmmDiff{} des
erreurs de marquage FO sont signalées par le mélange sur l'écart, \vFoStruct{} par le détecteur structurel.

\subsection{Là où ils sont utiles}
\begin{itemize}
\item \textbf{P-splines dans le modèle nul.} Si $\beta_0+\beta_1\ln T$ se révèle trop rigide sur données réelles, une
spline pénalisée~\cite{eilers1996} peut le remplacer comme partie lisse du modèle nul ; les tests de la
\cref{sec-detection} restent inchangés.
\item \textbf{Mélanges sur des sorties structurées.} Un mélange gaussien sur les sorties du détecteur ($\hat\beta_0$, $\hat\beta_1$,
les $z_t$, $Q$ et les rapports de dispersion) apprend les régimes habituels par région et par catégorie
dans un espace de faible dimension où chaque coordonnée a un sens.
\item \textbf{Région $\times$ catégorie} sert au mieux à regrouper la calibration de $\kappa$ et de l'échelle de niveau,
pas à former des groupes.
\item \textbf{Réseaux de neurones.} D'après la \cref{prop-nn}, sur les mécanismes que l'on sait nommer, un réseau ne peut
au mieux qu'apprendre le filtre adapté. Les réseaux trouvent leur place une fois qu'existe un historique d'incidents
étiquetés et des motifs que les gabarits ne couvrent pas ; ils doivent alors être entraînés sur les résidus du modèle
structurel et sur des injections synthétiques, pas sur les courbes brutes.
\end{itemize}

\section{Mise en œuvre et exploitation}\label{sec-ops}

\subsection{Package}
Le package Python \code{repo\_control} (numpy et scipy ; scikit-learn uniquement pour la référence par mélange gaussien)
implémente tous les objets de cette note :
\begin{center}\small
\begin{tabular}{p{2.9cm}p{12.2cm}}
\toprule
Module & Contenu\\
\midrule
\code{curves} & courbes zéro-coupon à forwards constants, dividendes discrets, forwards, bootstrap de l'\cref{alg-bootstrap}\\
\code{parity} & forward et facteur d'actualisation à partir des prix de consensus des calls et des puts (\cref{lem-parity})\\
\code{buckets} & périodes de dividendes, allocation, matrice d'allocation $A$ (\cref{lem-allocation})\\
\code{noise} & écart de panel, covariance de panels emboîtés, covariances des panels options et dividendes, plafonnement de la dispersion (\cref{sec-noise})\\
\code{pipeline} & instantanés Totem et FO $\to$ repo implicite, écart de portage, $\Sigma$, contexte de détection\\
\code{detect} & gabarits, filtres adaptés, calibrations familiale et de $\kappa$, attribution, réparation, test d'erreur de marquage FO\\
\code{control} & lignes de contrôle, jacobienne du bootstrap $J$, matérialité de la répartition (\cref{sec-control})\\
\code{pricer} & un pricer minimal qui produit les sensibilités par période pour les exemples\\
\code{synthetic}, \code{benchmark}, \code{baselines} & marché synthétique construit contributeur par contributeur et banc d'essai de la \cref{sec-benchmark}\\
\bottomrule
\end{tabular}
\end{center}
Chaque composant a un contrôle en forme fermée dans la suite de tests : aller-retour du bootstrap, parité, identité de
contrôle contre la réévaluation complète, formules d'écart de panel et de sortie de contributeurs, calibration du
détecteur sous l'hypothèse nulle et reconnaissance de chaque gabarit.

\subsection{Données de production}
\begin{center}\small
\begin{tabular}{p{2.4cm}p{10.0cm}p{2.4cm}}
\toprule
Donnée & Champs & Forme\\
\midrule
Options Totem & échéances ; strikes ; prix de consensus des calls et des puts ; contributeurs ; dispersion des forwards des contributeurs (ou une estimation) ; spot de référence & $K\times M$\\
Dividendes Totem & bornes des périodes ; niveau de consensus ; dispersion ; contributeurs (net des dividendes déjà versés) & $Y$\\
Marché FO & spot à l'heure de l'instantané Totem ; courbe des taux ; courbe de repo zéro-coupon aux échéances Totem ; prévision de dividendes & $K$ ; liste\\
Risque FO & sensibilités repo par pilier (projetées sur les échéances Totem) ; sensibilités dividendes par période Totem & $K+Y$ par portefeuille\\
Historique & douze à vingt-quatre instantanés mensuels et incidents connus & pour $\kappa$ et la validation\\
\bottomrule
\end{tabular}
\end{center}
Si la dispersion des forwards des contributeurs n'est pas publiée, on peut l'estimer à partir de la dispersion des calls et
des puts très dans la monnaie, ou la calibrer sur l'historique ; le détecteur a besoin de la bonne échelle par échéance, et
$\kappa$ corrige un biais commun.

\subsection{Procédure mensuelle}
\begin{enumerate}
\item Charger les instantanés Totem et FO à la même heure de valorisation ; vérifier le spot de référence.
\item Impliciter les forwards par parité ; examiner le facteur d'actualisation implicite, les résidus de parité et les
rapports de dispersion.
\item Allouer les périodes de dividendes Totem sur le profil de détachement FO ; impliciter le repo avec les dividendes
Totem et avec les dividendes FO.
\item Construire $\Sigma$ et lancer le détecteur avec le $\kappa$ et l'échelle de niveau de la région et de la catégorie.
\item Calculer les lignes de contrôle par portefeuille : lignes forward, dividende, non-forward et repo ; vérifier la
matérialité de la répartition.
\item Examiner les alertes avec leurs explications et les erreurs de marquage FO possibles ; utiliser la courbe réparée
lorsqu'une anomalie Totem est confirmée.
\item Consigner la décision prise pour chaque alerte ; ces étiquettes alimentent la calibration suivante.
\end{enumerate}

\subsection{Calibration et gouvernance}
$\kappa$ et l'échelle de niveau sont recalibrés chaque trimestre par région et par catégorie sur un historique sain, avec
un budget d'alertes de \vFPR{} des sous-jacents par mois. Le taux d'alerte réalisé et la part des explications acceptées
par les relecteurs sont suivis. Là où ils sont cotés, des ancres de marché complètent le détecteur : contrats à terme
cotés sur indices et sur dividendes, contrats à terme de rendement total qui cotent un spread de financement, et
commissions de prêt-emprunt de titres pour le court terme des actions.

\section{Conclusion}
Les courbes de repo implicite sont irrégulières parce qu'elles sont le résidu de deux moyennes de consensus calculées sur
des panels différents ; une partie de ce résidu ne peut être supprimée par aucune méthode qui ne voit que les données de
consensus (\cref{thm-nonid}). Deux conséquences en découlent. Pour une exposition linéaire en forward, le contrôle total
n'a pas besoin de la répartition, exactement en réévaluation complète et au premier ordre sur les lignes de contrôle
(\cref{thm-control,thm-exact}), et l'exposition non-forward $g$ indique, portefeuille par portefeuille, l'importance de
la répartition. Pour la détection, les mécanismes qui produisent réellement des courbes aberrantes ont des signatures
connues dans l'espace de portage ; des filtres adaptés, avec un modèle de bruit construit à partir des propres statistiques
de Totem, les détectent à un taux de fausses alertes calibré, les expliquent en unités naturelles et les distinguent des
véritables erreurs de marquage FO. L'étape suivante consiste à faire tourner la chaîne sur douze à vingt-quatre mois
d'instantanés Totem et FO réels, à calibrer les réglages par région et par catégorie, et à mesurer l'attribution sur des
incidents étiquetés.

\appendix
\section{Notations}\label{app-notation}
\begin{center}\small
\begin{tabular}{lp{12cm}}
\toprule
Symbole & Signification\\
\midrule
$S$, $r$, $R=rt$ & spot, taux zéro-coupon, taux cumulé\\
$b$, $B=bt$ & repo zéro-coupon, repo cumulé (portage)\\
$G=e^{R-B}$ & facteur de croissance du portage\\
$t_i$, $d_i$, $\tilde d_i=d_i/G(t_i)$ & dates de détachement, dividendes en numéraire, dividendes actualisés au portage\\
$D(T)$ & dividende cumulé actualisé au portage, \cref{eq-forward}\\
$T_k$, $k=1..K$ & échéances Totem (piliers)\\
$L_y$, $\pi_i$, $A_{ky}$ & niveau de dividendes de la période $y$, poids du profil, matrice d'allocation\\
$\Aset$, $\Bset$, $n_A$, $n_B$, $m$ & panels options et dividendes, leurs tailles, contributeurs communs aux deux\\
$\bA$, $\DA$, $\DB$ & moyennes des vues de repo et de dividendes du panel options ; consensus du panel dividendes\\
$y_k$ & écart de portage Totem moins FO, \cref{eq-y}\\
$X=[T,T\ln T]$, $\beta$ & modèle nul (écart de repo lisse) et ses coefficients\\
$\Sigma$, $\kappa$ & covariance du bruit de consensus, échelle de variance\\
$t$, $\tb$, $z_t$, $Q$ & gabarit, sa projection blanchie, statistique de filtre adapté, chi-deux global\\
$s_b$, $s_D$, $J$, $g$ & sensibilités FO au repo et aux dividendes, jacobienne du repo implicite, exposition non-forward\\
\bottomrule
\end{tabular}
\end{center}

\section{Démonstrations et remarques complémentaires}\label{app-proofs}

\begin{lemma}[Actualisations hétérogènes]
Si le contributeur $j$ actualise avec $DF_j$, le zéro de $\overline{C}-\overline{P}$ vaut $K^\ast=\sum_jDF_jF_j/\sum_jDF_j$,
une moyenne des forwards des contributeurs pondérée par l'actualisation. Son écart à la moyenne simple vaut
$\Cov(DF_j,F_j)/\overline{DF}$, du second ordre.
\end{lemma}
\begin{proof}
$\overline{C}-\overline{P}=\frac1n\sum_jDF_j(F_j-K)$ s'annule en $K^\ast$. Ensuite
$K^\ast-\overline F=\bigl(\overline{DF\,F}-\overline{DF}\,\overline F\bigr)/\overline{DF}$.
\end{proof}

\begin{lemma}[Rapport signal sur bruit exact dans un plan apparié]
Si un scénario ajoute un effet $\eta$ à l'écart de portage d'un sous-jacent sans modifier son bruit de consensus (même
graine), alors $y^{\text{scénario}}-y^{\text{sain}}=\eta$ et le rapport signal sur bruit $\|PL^{-1}\eta\|$ du
\cref{cor-power} est connu exactement, sans bruit. Pour un seul gabarit ($\eta=a\,t$), il vaut $|a|\,\|\tb\|$.
\end{lemma}
\begin{proof}
Le bruit s'élimine dans la différence ; pour $\eta=at$, $PL^{-1}\eta=a\tb$.
\end{proof}

\begin{lemma}[Calibration de Student du test d'erreur de marquage FO]
Soit $\hat\beta_0\sim\N(\beta_0,c\,\sigma^2)$ et $\hat\sigma^2$ une estimation indépendante avec
$\nu\hat\sigma^2/\sigma^2\sim\chi^2_\nu$. Alors $(\hat\beta_0-\beta_0)/\sqrt{c\hat\sigma^2}\sim t_\nu$.
\end{lemma}
\begin{proof}
C'est la définition de la loi de Student : une normale centrée réduite divisée par la racine d'un $\chi^2_\nu/\nu$
indépendant.
\end{proof}
Avec des dispersions estimées sur $n_A$ contributeurs, $\nu\approx n_A-1$ sur le court terme, là où se lit l'incertitude
du niveau de repo ; l'échelle de niveau $s$ corrige ensuite ce que cette approximation laisse échapper.

\begin{remark}[Terme du second ordre sur plusieurs piliers]
La \cref{prop-second} s'étend pilier par pilier : pour $V=\Phi(F_1,\dots,F_K)$ avec des échéances aux piliers, la dépendance
du total au premier ordre vis-à-vis du panel dividendes vaut
$-\sum_k\partial_k\Phi\,F_k\bigl[\tfrac12T_k^2\Delta b_k^2+T_k\Delta b_k\Delta D_k/(S-D_k)\bigr]$, plus les termes croisés créés
par l'interpolation du portage entre piliers, tous du second ordre.
\end{remark}

\section{Spécification du marché synthétique}\label{app-synthetic}
\begin{itemize}
\item \textbf{Dividendes.} Actions UE : dividende annuel détaché entre le 20 avril et le 30 mai ($35\,\%$ des titres versent
aussi un acompte de $40\,\%$ en novembre) ; actions US : trimestriel ; titres difficiles à emprunter : annuel ou
trimestriel ; indice UE : flux mensuel avec des poids $(1,2,4,14,24,18,6,4,3,4,3,2)$ par mois civil. Rendement et
croissance tirés par sous-jacent dans les intervalles du tableau des catégories.
\item \textbf{Repo.} $b(T)=b_\infty+(b_0-b_\infty)e^{-T/\tau}$, avec $b_0$, $b_\infty$ et $\tau$ tirés par sous-jacent.
\item \textbf{Vues dividendes des contributeurs.} $d_{ij}=d_i\bigl(1+\sigma(t_i)[0{,}7\,z_j+\sqrt{1-0{,}7^2}\,e_{ij}]\bigr)$, avec
$\sigma(t)$ affine en $t$ (propre à chaque catégorie), $z_j$ persistant et $e_{ij}$ idiosyncratique ; les dividendes annoncés
(avant $0{,}15$ an) ont un dixième de la dispersion.
\item \textbf{Vues repo des contributeurs.} $b_j(T)=b(T)+\sigma_b\bigl(u_j+0{,}6\,v_j\,\tau(T)\bigr)$, où $\tau(T)$ est une
inclinaison linéaire en $\ln T$ de $-0{,}5$ à $0{,}5$ ; $\sigma_b$ vaut $15$, $10$, $60$ et $8$\,pb pour les quatre catégories.
\item \textbf{Options.} Strikes de $80$ à $120\,\%$ du spot ; volatilité avec structure par terme et smile de $-0{,}10$ par
unité de log-moneyness, bruit de volatilité d'un point par contributeur, bruit de prix de $2\times10^{-5}$ du spot.
\item \textbf{Consensus.} Moyenne sur le panel après exclusion des cotations à plus de quatre écarts absolus médians
normalisés de la médiane.
\item \textbf{FO.} Dividendes avec un bruit de $1\,\%$, repo avec une erreur de niveau de $3$\,pb ; spot égal au spot de
référence Totem.
\end{itemize}

\section{Reproductibilité}\label{app-repro}
Tous les nombres, tableaux et figures de cette note sont produits par le dépôt :
\begin{center}\small
\begin{tabular}{l}
\code{pip install -e .[benchmark,report,dev]}\\
\code{pytest}\\
\code{python scripts/run\_benchmark.py} \quad (environ \vRuntime{} minutes)\\
\code{PYTHONPATH=. python scripts/build\_latex.py}
\end{tabular}
\end{center}

\begin{thebibliography}{9}
\bibitem{eilers1996} P.~H.~C. Eilers et B.~D. Marx. Flexible smoothing with B-splines and penalties.
\emph{Statistical Science}, 11(2):89--121, 1996.
\bibitem{hull} J.~C. Hull. \emph{Options, Futures, and Other Derivatives}. Pearson, toute édition récente.
\bibitem{kay1998} S.~M. Kay. \emph{Fundamentals of Statistical Signal Processing, Volume II: Detection Theory}.
Prentice Hall, 1998.
\bibitem{lovell1963} M.~C. Lovell. Seasonal adjustment of economic time series and multiple regression analysis.
\emph{Journal of the American Statistical Association}, 58(304):993--1010, 1963.
\bibitem{mallat1993} S.~G. Mallat et Z.~Zhang. Matching pursuits with time-frequency dictionaries.
\emph{IEEE Transactions on Signal Processing}, 41(12):3397--3415, 1993.
\bibitem{eu2016} Règlement délégué (UE) 2016/101 de la Commission du 26 octobre 2015 complétant le règlement (UE)
n\textsuperscript{o}~575/2013 par des normes techniques de réglementation relatives à l'évaluation prudente.
\bibitem{scharf1994} L.~L. Scharf et B.~Friedlander. Matched subspace detectors. \emph{IEEE Transactions on Signal
Processing}, 42(8):2146--2157, 1994.
\bibitem{sidak1967} Z.~Šidák. Rectangular confidence regions for the means of multivariate normal distributions.
\emph{Journal of the American Statistical Association}, 62(318):626--633, 1967.
\end{thebibliography}

\end{document}
